\batchmode
\makeatletter
\def\input@path{{"C:/Users/tkpme/Dropbox/BookPCRA/Springer Production March 2026/Ch Appendix A Math for PCRA/"}}
\makeatother
\documentclass[graybox,envcountchap,sectrefs]{svmono}\usepackage[]{graphicx}\usepackage[]{xcolor}
% maxwidth is the original width if it is less than linewidth
% otherwise use linewidth (to make sure the graphics do not exceed the margin)
\makeatletter
\def\maxwidth{ %
  \ifdim\Gin@nat@width>\linewidth
    \linewidth
  \else
    \Gin@nat@width
  \fi
}
\makeatother

\definecolor{fgcolor}{rgb}{0.345, 0.345, 0.345}
\newcommand{\hlnum}[1]{\textcolor[rgb]{0.686,0.059,0.569}{#1}}%
\newcommand{\hlsng}[1]{\textcolor[rgb]{0.192,0.494,0.8}{#1}}%
\newcommand{\hlcom}[1]{\textcolor[rgb]{0.678,0.584,0.686}{\textit{#1}}}%
\newcommand{\hlopt}[1]{\textcolor[rgb]{0,0,0}{#1}}%
\newcommand{\hldef}[1]{\textcolor[rgb]{0.345,0.345,0.345}{#1}}%
\newcommand{\hlkwa}[1]{\textcolor[rgb]{0.161,0.373,0.58}{\textbf{#1}}}%
\newcommand{\hlkwb}[1]{\textcolor[rgb]{0.69,0.353,0.396}{#1}}%
\newcommand{\hlkwc}[1]{\textcolor[rgb]{0.333,0.667,0.333}{#1}}%
\newcommand{\hlkwd}[1]{\textcolor[rgb]{0.737,0.353,0.396}{\textbf{#1}}}%
\let\hlipl\hlkwb

\usepackage{framed}
\makeatletter
\newenvironment{kframe}{%
 \def\at@end@of@kframe{}%
 \ifinner\ifhmode%
  \def\at@end@of@kframe{\end{minipage}}%
  \begin{minipage}{\columnwidth}%
 \fi\fi%
 \def\FrameCommand##1{\hskip\@totalleftmargin \hskip-\fboxsep
 \colorbox{shadecolor}{##1}\hskip-\fboxsep
     % There is no \\@totalrightmargin, so:
     \hskip-\linewidth \hskip-\@totalleftmargin \hskip\columnwidth}%
 \MakeFramed {\advance\hsize-\width
   \@totalleftmargin\z@ \linewidth\hsize
   \@setminipage}}%
 {\par\unskip\endMakeFramed%
 \at@end@of@kframe}
\makeatother

\definecolor{shadecolor}{rgb}{.97, .97, .97}
\definecolor{messagecolor}{rgb}{0, 0, 0}
\definecolor{warningcolor}{rgb}{1, 0, 1}
\definecolor{errorcolor}{rgb}{1, 0, 0}
\newenvironment{knitrout}{}{} % an empty environment to be redefined in TeX

\usepackage{alltt}
\usepackage{mathptmx}
\usepackage{helvet}
\usepackage{courier}
\usepackage[T1]{fontenc}
\usepackage[utf8]{inputenc}
\usepackage{geometry}
\geometry{verbose,tmargin=1in,bmargin=1in,lmargin=1in,rmargin=1in}
\setcounter{tocdepth}{3}
\setlength{\parskip}{\medskipamount}
\setlength{\parindent}{0pt}
\usepackage{float}
\usepackage{units}
\usepackage{textcomp}
\usepackage{url}
\usepackage{amsmath}
\usepackage{amssymb}
\usepackage[unicode=true,pdfusetitle,
 bookmarks=true,bookmarksnumbered=true,bookmarksopen=true,bookmarksopenlevel=1,
 breaklinks=false,pdfborder={0 0 1},backref=false,colorlinks=false]
 {hyperref}
\hypersetup{
 pdfstartview=XYZ null null 1.25,pdfpagemode=UseNone}

\makeatletter

%%%%%%%%%%%%%%%%%%%%%%%%%%%%%% LyX specific LaTeX commands.
\newcommand{\lyxmathsym}[1]{\ifmmode\begingroup\def\b@ld{bold}
  \text{\ifx\math@version\b@ld\bfseries\fi#1}\endgroup\else#1\fi}


%%%%%%%%%%%%%%%%%%%%%%%%%%%%%% User specified LaTeX commands.
%%%%%%%%%%%%%%%%%%%% book.tex %%%%%%%%%%%%%%%%%%%%%%%%%%%%%
%
% sample root file for the chapters of your "monograph"
%
% Use this file as a template for your own input.
%
%%%%%%%%%%%%%%%% Springer-Verlag %%%%%%%%%%%%%%%%%%%%%%%%%%


% RECOMMENDED %%%%%%%%%%%%%%%%%%%%%%%%%%%%%%%%%%%%%%%%%%%%%%%%%%%


% choose options for [] as required from the list
% in the Reference Guide

\usepackage{makeidx}% allows index generation
% standard LaTeX graphics tool
                             % when including figure files
\usepackage{multicol}% used for the two-column index
\usepackage[bottom]{footmisc}% places footnotes at page bottom

\usepackage{hyperref}
\def\UrlBreaks{\do\/\do-}

% see the list of further useful packages
% in the Reference Guide

\makeindex             % used for the subject index
                       % please use the style svind.ist with
                       % your makeindex program


%\usepackage[ style=authoryear,natbib=true,giveninits=true,backend=biber,maxcitenames=2]{biblatex}
\usepackage[style=authoryear,natbib=true,firstinits=false,backend=biber,maxcitenames=2]{biblatex}


\setlength{\bibitemsep}{0.5\baselineskip}
\DeclareNameAlias{sortname}{family-given}

\IfFileExists{C:/Users/Doug/Dropbox/BookPCRA/mpslbookCombined.bib}
{\addbibresource[location=local]{C:/Users/Doug/Dropbox/BookPCRA/mpslbookCombined.bib}}
{   \IfFileExists{C:/Users/tkpme/Dropbox/BookPCRA/mpslbookCombined.bib}
    {\addbibresource[location=local]{C:/Users/tkpme/Dropbox/BookPCRA/mpslbookCombined.bib}}
    {   \IfFileExists{C:/Users/stoya/Dropbox/BookPCRA/mpslbookCombined.bib}
        {\addbibresource[location=local]{C:/Users/stoya/Dropbox/BookPCRA/mpslbookCombined.bib}}
        {   \IfFileExists{C:/Users/4thauthor/Dropbox/BookPCRA/mpslbookCombined.bib}
            {\addbibresource[location=local]{C:/Users/4thauthor/Dropbox/BookPCRA/mpslbookCombined.bib}}
            {}
        }
    }
}
\renewcommand\cite{\citet}

\usepackage{amsmath}
\usepackage{amssymb}
\usepackage[p,osf]{stickstootext}
\usepackage[stix2,vvarbb]{newtxmath}
\usepackage[varqu,varl]{inconsolata}
% \usepackage{txfonts}
\usepackage{upgreek}

\DeclareMathOperator{\argmax}{arg\,max}
\DeclareMathOperator{\argmin}{arg\,min}
\usepackage{array}
\usepackage{arydshln}
\usepackage{booktabs}
\usepackage{threeparttablex}
\usepackage{appendix}

\usepackage[utf8]{inputenc}
%%%%%%%%%%%%%%%%%%%%%%%%%%%%%%%%%%%%%%%%%%%
\IfFileExists{C:/Users/Doug/Dropbox/BookPCRA/macroSP_August_05_2021.tex}
{\input{C:/Users/Doug/Dropbox/BookPCRA/macroSP_August_05_2021.tex}}
{   \IfFileExists{C:/Users/tkpme/Dropbox/BookPCRA/macroSP_August_05_2021.tex}
    {\input{C:/Users/tkpme/Dropbox/BookPCRA/macroSP_August_05_2021.tex}}
    {   \IfFileExists{C:/Users/stoya/Dropbox/BookPCRA/macroSP_August_05_2021.tex}
        {\input{C:/Users/stoya/Dropbox/BookPCRA/macroSP_August_05_2021.tex}}
        {   \IfFileExists{C:/Users/4thAuthor/Dropbox/BookPCRA/macroSP_August_05_2021.tex}
            {\input{C:/Users/4thAuthor/Dropbox/BookPCRA/macroSP_August_05_2021.tex}}
            {}
        }
    }
}

\belowsink=90pt

%%%%%%%%%%%%%%%%%%%%%%%%%%%%%%%%%%%%%%%%%%%%%

\makeatother
\IfFileExists{upquote.sty}{\usepackage{upquote}}{}
\begin{document}
\begin{appendices}


\setcounter{chapter}{0} 

\chapter{MATHEMATICAL TOOLS FOR PCRA\label{chap:AppendixA}}

\vspace{0.3in}

\textbf{Planned for publication in 2026 by R. Douglas Martin, Thomas
K. Philips, Bernd Scherer and Kirk Li. All rights reserved. © Copyright
2025}

\vspace{0.3in}

\begin{center}
10/12/25
\par\end{center}

\tocchpnum=30pt 
\tocsecnum=33pt 
\tocsubsecnum=40pt 
\tocsubsubsecnum=48pt 
\tocparanum=55pt 
\tocsubparanum=63pt
\calctocindent

\tableofcontents{}\newpage{}

This Appendix contains a variety of basic mathematical results that
prove useful in the context of portfolio construction and risk management.
We expect that much of the material provided here will be used primarily
to refresh the reader's memory, thereby facilitating its use throughout
this book. Our treatment of matrices and Linear Algebra owes a great
deal to \citet{Gentle2017} and \citet{Strang2016}, and we strongly
recommend both books to anyone with an interest in this under rated
field. In addition, we refer readers who have an interest in optimization
to \citet{LuenbergerYe2008}.

\section{Vectors and Vector Spaces}

\subsection{Vectors}

A \emph{vector} is a collection of objects that supports two arithmetic
operations: vector addition (denoted +) and scalar multiplication
(denoted $\cdot$). It is useful to think of a vector as a list of
$n$ real numbers that correspond to a point in $n$--dimensional
Euclidean space with $n$ orthogonal axes, as this interpretation
clearly supports elementwise addition provided both vectors have the
same number of elements, i.e., if

\begin{equation}
\mathrm{\mathbf{u}}=\left[\begin{array}{c}
u_{1}\\
\vdots\\
u_{n}
\end{array}\right],\ \mathrm{\mathbf{v}}=\left[\begin{array}{c}
v_{1}\\
\vdots\\
v_{n}
\end{array}\right]\label{eq:Vector1}
\end{equation}

then

\begin{align}
\mathrm{\mathbf{u}}\,+\,\mathrm{\mathbf{v}} & =\left[\begin{array}{c}
u_{1}\\
\vdots\\
u_{n}
\end{array}\right]\,+\,\left[\begin{array}{c}
v_{1}\\
\vdots\\
v_{n}
\end{array}\right]\label{eq:Vector2}\\[0.5em]
 & =\left[\begin{array}{c}
u_{1}\,+\,v_{1}\\
\vdots\\
u_{n}\,+\,v_{n}
\end{array}\right].\label{eq:Vector3}
\end{align}

In addition, it supports scalar multiplication according to

\begin{equation}
c\,\cdot\,\mathrm{\mathbf{u}}=\left[\begin{array}{c}
c\,u_{1}\\
\vdots\\
c\,u_{n}
\end{array}\right].\label{eq:Vector4}
\end{equation}

However, not all vectors are defined in this way. In particular, polynomials
and square--integrable functions of one real variable are also vectors,
as both support vector addition\footnote{The sum of two polynomials is also a polynomial, and the same is true
for square--integrable functions of one real variable.} and scalar multiplication. 

One type of vector that we repeatedly encounter in this book is a
vector of \emph{security returns}, each element of which is lower
bounded by $-100\%$. Another is a vector of \emph{security weights}
associated with a portfolio. These weights may be \emph{constrained}
(e.g. the long--only constraint requires all weights to be positive
and to sum to 1) or \emph{unconstrained} (i.e. the weights may be
positive or negative, and may not sum to 1, as is often the case when
a portfolio is levered).

Vectors also support vector multiplication via the \emph{dot product,
}or \emph{inner product}\footnote{The term dot product is typically used when the vector space under
consideration is $n$--dimensional Euclidean space. In more general
vector spaces, the term inner product is more commonly used. There
is another type of vector product, the \emph{cross product}, which
sees application mainly in physics. We ignore it here.}. The dot product of two vectors of equal length, $\mathrm{\mathbf{u}}$
and $\mathrm{\mathbf{v}}$, is defined to be\footnote{This definition of an inner product can be extended to cover complex
numbers as well, but we restrict our attention to real numbers as
complex numbers are exceedingly rare in portfolio construction.}

\begin{equation}
\mathrm{\mathbf{u}}\,\cdotp\,\mathrm{\mathbf{v}}=\sum_{i\,=\,1}^{n}u_{i}\,v_{i}.\label{eq:DotProduct}
\end{equation}

More generally, the inner product of $\mathrm{\mathbf{u}}$ and $\mathrm{\mathbf{v}}$
is written as $\left\langle \mathrm{\mathbf{u}},\,\mathrm{\mathbf{v}}\right\rangle $,
and is a real number that satisfies the following criteria:
\begin{enumerate}
\item Non--negativity: $\mathrm{\mathbf{\left\langle \mathrm{\mathbf{u}},\,\mathrm{\mathbf{0}}\right\rangle =}\,0\ \forall\,\mathrm{\mathbf{u}}},$
and $\left\langle \mathrm{\mathbf{u}},\,\mathrm{\mathbf{u}}\right\rangle >0$
if $\mathrm{\mathbf{u}}\ne\mathbf{0}$
\item Commutativity: $\left\langle \mathrm{\mathbf{u}},\,\mathrm{\mathbf{v}}\right\rangle =\left\langle \mathrm{\mathbf{v}},\,\mathrm{\mathbf{u}}\right\rangle $
\item Scalar multiplication: $\left\langle c\,\mathrm{\mathbf{u}},\,\mathrm{\mathbf{v}}\right\rangle =c\,\left\langle \mathrm{\mathbf{u}},\,\mathrm{\mathbf{v}}\right\rangle $
for all scalars $c$
\item Distributivity: $\left\langle \mathrm{\mathbf{u}\,+\,}\mathbf{w},\,\mathrm{\mathbf{v}}\right\rangle =\left\langle \mathrm{\mathbf{u}},\,\mathrm{\mathbf{v}}\right\rangle \,+\,\left\langle \mathrm{\mathbf{w}},\,\mathrm{\mathbf{v}}\right\rangle $
\end{enumerate}
Two vectors, $\mathrm{\mathbf{u}}$ and $\mathrm{\mathbf{v}}$, are
said to be \emph{orthogonal} if 

\begin{equation}
\left\langle \mathrm{\mathbf{u}},\,\mathrm{\mathbf{v}}\right\rangle =0.\label{eq:OrthogonalVectors}
\end{equation}

The \emph{norm} of a vector is the square root of its inner product
with itself.

\begin{equation}
\left\Vert \mathrm{\mathbf{u}}\right\Vert =\sqrt{\left\langle \mathrm{\mathbf{u}},\,\mathrm{\mathbf{u}}\right\rangle }.\label{eq:NormU}
\end{equation}

In Euclidean space, the norm of a vector is equal to its length.

\subsection{The Distance Between Two Vectors}

The vector

\begin{equation}
\mathrm{\mathbf{u}}=\left[\begin{array}{r}
1\\
-2
\end{array}\right]\label{eq:VectorWith2Elements}
\end{equation}

corresponds to the point $x=1,\,y=-2$ in the $X-Y$ plane. We can
apply Pythagoras' Theorem to compute its norm, or the distance from
the origin to the point $\left(1,\,-2\right)$:

\begin{align}
\left\Vert \mathrm{\mathbf{u}}\right\Vert  & =\sqrt{1^{2}\,+\,\left(-2\right)^{2}}\nonumber \\
 & =\sqrt{5}.\label{eq:LengthOfVectorWith2Elements}
\end{align}

In general, the norm of a vector with $n$ elements in Euclidean space
is

\begin{equation}
\left\Vert \mathrm{\mathbf{u}}\right\Vert =\sqrt{u_{1}^{2}\,+\,\cdots\,+\,u_{n}^{2}},\label{eq:LengthOfVector}
\end{equation}

and this concept is easily generalized to obtain the \emph{distance}
between two vectors:

\begin{align}
d_{\mathrm{\mathbf{u}},\,\mathrm{\mathbf{v}}} & \triangleq\left\Vert \mathrm{\mathbf{u}}-\mathrm{\mathbf{v}}\right\Vert \nonumber \\
 & =\sqrt{\left(u_{1}-v_{1}\right)^{2}\,+\,\cdots\,+\,\left(u_{n}-v_{n}\right)^{2}}.\label{eq:DistanceBetweenVectors}
\end{align}

It is important to note that this definition of distance is not unique:
the distance between two square--integrable functions of one real
variable, $f(x)$ and $g(x)$, for example, can be written as

\begin{equation}
d_{f,\,g}=\sqrt{\int\left(f(x)\,-\,g(x)\right)^{2}dx}.\label{eq:DistanceBetweenTwoFunctions}
\end{equation}

Regardless of how it is defined, a valid measure of distance must
satisfy three properties:
\begin{enumerate}
\item Non-negativity: For all vectors $\mathrm{\mathbf{u}},\,\mathrm{\mathbf{v}},\ d_{\mathrm{\mathbf{u}},\mathrm{\,\mathbf{v}}}\geqslant0.$
\item Symmetry: For all vectors $\mathrm{\mathbf{u}},\,\mathrm{\mathbf{v}},\ d_{\mathrm{\mathbf{u}},\,\mathrm{\mathbf{v}}}=d_{\mathrm{\mathbf{v}},\,\mathrm{\mathbf{u}}}$.
\item Triangle Inequality: For all vectors $\mathrm{\mathbf{u}},\,\mathrm{\mathbf{v}},\,\mathrm{\mathbf{w}},\ d_{\mathrm{\mathbf{u}},\,\mathrm{\mathbf{w}}}\leqslant d_{\mathrm{\mathbf{u}},\,\mathrm{\mathbf{v}}}\,+\,d_{\mathrm{\mathbf{v}},\,\mathrm{\mathbf{w}}}.$ 
\end{enumerate}
The triangle inequality get its name from its geometric interpretation:
the sum of the lengths of any two sides of a triangle must be at least
as large as its third side. When the vectors in question are lists
of real numbers, and the measure of distance between vectors is the
Euclidean distance, the triangle inequality is equivalent to the Cauchy--Schwarz
inequality for real numbers:

\begin{equation}
\left(\sum_{i\,=\,1}^{n}u_{i}\,v_{i}\right)^{2}\leqslant\sum_{i\,=\,1}^{n}u_{i}^{2}\,\cdotp\,\sum_{i\,=\,1}^{n}v_{i}^{2},\label{eq:CauchySchwarz}
\end{equation}
with equality if and only if $v_{i}=c\,u_{i}$ for some constant $c$.

While proofs of these inequalities abound, we recommend reviewing
their treatment in \citet{Steele2004}, which is both comprehensive
and accessible to the non--specialist.  The dot product admits an
equivalent, geometric, definition:

\begin{equation}
\mathrm{\mathbf{u}}\,\cdotp\,\mathrm{\mathbf{v}}=\left\Vert \mathrm{\mathbf{u}}\right\Vert \,\left\Vert \mathrm{\mathbf{v}}\right\Vert \,\cos\theta,\label{eq:DotProduct2}
\end{equation}
where $\theta$ is the angle between $\mathrm{\mathbf{u}}$ and $\mathrm{\mathbf{v}}$.
Orthogonality is equivalent to condition that $\theta=90^{\lyxmathsym{\textdegree}}$.

\subsection{Vector Spaces}

A \emph{vector space} adds structure to sets of vectors and the arithmetic
operations associated with them. A vector space $\mathcal{{S}}$ has
the following properties:
\begin{enumerate}
\item Additive Closure: For all vectors $\mathrm{\mathbf{u}},\,\mathrm{\mathbf{v}}\in\mathcal{{S}},\ \mathrm{\mathbf{u}}\,+\,\mathrm{\mathbf{v}}\in\mathcal{{S}}$.
\item Multiplicative Closure: For all vectors $\mathrm{\mathbf{u}}\in\mathcal{{S}},$
and for all scalars $c,\ c\,\cdot\,\mathbf{u}\in\mathcal{{S}}$.
\item Additive commutativity: For all vectors $\mathrm{\mathbf{u}},\,\mathrm{\mathbf{v}}\in\mathcal{{S}},\ \mathrm{\mathbf{u}}\,+\,\mathrm{\mathbf{v}}=\mathrm{\mathbf{v}}\,+\,\mathrm{\mathbf{u}}$
. 
\item Distributivity with respect to scalar and vector addition: For all
vectors $\mathrm{\mathbf{u}},\,\mathrm{\mathbf{v}}\in\mathcal{{S}}$,
and for all scalars $c,\,d,\ (c\,+\,d)\cdot\left(\mathrm{\mathbf{u}}\,+\,\mathrm{\mathbf{v}}\right)=\mathit{c}\cdot\mathbf{\mathrm{\mathbf{u}}}\,+\,\mathit{c}\,\cdot\,\mathrm{\mathbf{v}}\,+\,\mathit{d}\,\cdot\,\mathrm{\mathbf{u}}\,+\,\mathit{d}\,\cdot\,\mathrm{\mathbf{v}}\in\mathcal{{S}}$.
\item Additive Associativity: For all vectors $\mathrm{\mathbf{u}},\,\mathrm{\mathbf{v},}\,\mathrm{\mathbf{w}}\in\mathcal{{S}},\ \left(\mathrm{\mathbf{u}}\,+\,\mathrm{\mathbf{v}}\right)\,+\,\mathrm{\mathbf{w}}=\mathrm{\mathbf{u}}\,+\,\left(\mathrm{\mathbf{v}}\,+\,\mathrm{\mathbf{w}}\right).$
\item Multiplicative Associativity: For all vectors $\mathrm{\mathbf{u}}\in\mathcal{{S}},$
and for all scalars $c,\,d$, $\mathrm{\ }\left(c\,d\right)\,\cdot\,\mathrm{\mathbf{w}}=c\,\cdot\,\left(d\,\cdot\,\mathrm{\mathbf{w}}\right).$
\item Existence of additive and multiplicative identities. There exists
a vector $\mathbf{\mathrm{\mathbf{0}}}\in\mathcal{{S}}$ and a scalar
1, such that for all $\mathrm{\mathbf{u}}\in\mathcal{{S}},\ \mathrm{\mathbf{u}\mathrm{\,+\,\mathbf{0}}}=\mathrm{\mathbf{u}}$
and $1\,\cdot\,\mathbf{u}=\mathbf{u}$.
\item Existence of an additive inverse. For all vectors, $\mathrm{\mathbf{u}}\in\mathcal{{S}},$
there exists a vector $\mathrm{\mathbf{-u}}\in\mathcal{{S}}$ such
that $\mathrm{\mathbf{u}}\,+\,\mathrm{\mathbf{-u}}=\mathrm{\mathbf{0}}$.
\end{enumerate}
It is important to check that each of these properties is satisfied
before declaring a set of vectors to be a vector space. For example,
if $\mathrm{\mathbf{u}}$ is a vector of security returns, and if
$c$ is a real number, and if vector addition and scalar multiplication
are defined elementwise and follow the normal rules of arithmetic,
$c\,\cdotp\,\mathrm{\mathbf{u}}$ does not always satisfy a basic
property of returns: they cannot be smaller than $-100\%$. If, on
the other hand, $\mathrm{\mathbf{u}}$ is a vector of \emph{logarithmic}
\emph{returns} (see \AssetReturnsChapter\  for a discussion of various
types of returns) it is easy to see that $\mathcal{{S}}$ is a vector
space.

An \emph{inner product space} is a vector space in which the inner
product is well defined.

A \emph{subspace} of a vector space is a subset of $\mathcal{{S}}$
that satisfies all the requirements of a vector space. In particular,
vector addition and scalar multiplication must always result in a
vector that lies in the subspace. 

For example, let $\mathcal{{S}}_{XY}$ be the set of all vectors with
2 real components, i.e. vectors in the $X\lyxmathsym{–}Y$ plane,
with elementwise vector addition and scalar multiplication. $\mathcal{{S}}_{Y}$,
the subset of vectors whose first component is 0 (i.e. vectors that
lie on the $Y$ axis), is a subspace of $\mathcal{{S}}_{XY}$, as
is $\mathcal{{S}}_{X}$, the subset of vectors whose second component
is 0 (i.e. vectors that lie on the $X$ axis). 

Formally, a set of vectors $\mathrm{\mathbf{v}}_{1,},\mathrm{\mathbf{v}}_{2},\cdots,\mathrm{\mathbf{v}}_{k}$
which are not all equal to $\mathbf{0}$ is said to be linearly independent
if

\begin{equation}
c_{1}\,\mathrm{\mathbf{v}}_{1}\,+\,c_{2}\,\mathrm{\mathbf{v}}_{2}\,+\,\cdots\,+\,c_{k}\,\mathrm{\mathbf{v}}_{k}=\mathbf{0}\Longleftrightarrow c_{1}\,=\,c_{2}\,=\,\cdots\,=\,c_{k}\,=\,0.\label{eq:VectorSpaceDimension}
\end{equation}

The \emph{dimension} of a vector space is the maximum number of linearly
independent vectors contained in it, and the dimension of a subspace
cannot exceed that of its parent space. 

A set of linearly independent vectors $\mathcal{{B}}$ is said to
be a \emph{basis} for $\mathcal{{S}}$, or to \emph{span} $\mathcal{{S}}$,
if any vector $\mathrm{\mathbf{v}}\in\mathcal{{S}}$ can be written
as a unique linear combination of the vectors in $\mathcal{{B}}$.
Bases are not unique, though one basis may prove more convenient than
another in a particular application. All bases must have the same
number of elements, and this must equal the dimension of the vector
space. 

If the vectors in a basis are all orthogonal to each other, the basis
is called an \emph{orthogonal basis}. If, in addition, the length
of each basis vector is 1, it is said to be an \emph{orthonormal basis}.

It is easy to see why $\mathcal{{S}}_{XY}$ is a two dimensional vector
space: there are no non--zero values of $c_{1}$ and $c_{2}$ such
that

\begin{equation}
c_{1}\left[\begin{array}{c}
0\\
1
\end{array}\right]\,+\,c_{2}\left[\begin{array}{c}
1\\
0
\end{array}\right]=\left[\begin{array}{c}
0\\
0
\end{array}\right],\label{eq:2DVectorSpace}
\end{equation}

and every vector in the $X-Y$ plane can be represented as a unique
combination of $c_{1}$ and $c_{2}$. It follows that 

\begin{equation}
{B}=\left\{ \left[\begin{array}{c}
0\\
1
\end{array}\right],\ \left[\begin{array}{c}
1\\
0
\end{array}\right]\right\} \label{eq:BasisForXYplane}
\end{equation}
 is an orthonormal basis for $\mathcal{{S}}_{XY}$\footnote{Likewise, the monomials $\left\{ 1,\,x,\,x^{2,},\cdots\right\} $form
a basis for the vector space of polynomials.}.

Likewise, it is easy to see why $\mathcal{{S}}_{Y}$ is one--dimensional:
for any two vectors whose first component is 0, 

\begin{equation}
c_{1}\left[\begin{array}{c}
0\\
a
\end{array}\right]\,-\,\frac{ac_{1}}{b}\,\left[\begin{array}{c}
0\\
b
\end{array}\right]=\left[\begin{array}{c}
0\\
0
\end{array}\right],\label{eq:1DVectorSpace}
\end{equation}

provided that $a$ and $b$ are not 0. Furthermore, any point on the
$Y$ axis can be represented as a multiple of either vector. Exchanging
the first and the second components of both vectors makes clear that
the same argument applies to ${S}_{X}$.

\subsection{Projections of a Vector Space}

Exact, but computationally demanding, calculations using vectors in
a high dimensional vector space $\mathcal{{S}}_{1}$ are often approximated
using representations of the vectors in a lower dimensional subspace
$\mathcal{{S}}_{2}$. If $\mathcal{{B}}_{2}$ is an orthogonal basis
for $\mathcal{{S}}_{2}$, the best approximate representation of a
vector $\mathrm{\mathbf{u}_{1}}\in\mathcal{{S}}_{1}$ in $\mathcal{{S}}_{2}$
is its \emph{projection}, $\mathrm{\mathbf{p}}_{1}$, which is given
by
\begin{equation}
\mathrm{\mathbf{p}}_{1}=\sum_{\mathrm{\mathbf{b}}_{i}\in\mathcal{{B}}_{2}}\dfrac{\left\langle \mathrm{\mathbf{u}}_{1},\,\mathrm{\mathbf{b}}_{i}\right\rangle }{\left\Vert \mathrm{\mathbf{b}}_{i}\right\Vert ^{2}}\mathrm{\mathbf{b}}_{i}\label{eq:Projection}
\end{equation}

In essence, the projection first captures, and then aggregates, all
that $\mathrm{\mathbf{u}}_{1}$ has in common with each of the orthogonal
basis vectors of $\mathcal{{S}}_{2}$. The residual 

\begin{align}
\mathbf{r}_{1} & =\mathbf{u}_{1}\,-\,\mathbf{p}_{1}\nonumber \\
 & =\mathbf{u}_{1}\,-\,\sum_{\mathrm{\mathbf{b}}_{i}\in\mathcal{{B}}_{2}}\dfrac{\left\langle \mathrm{\mathbf{u}}_{1},\,\mathrm{\mathbf{b}}_{i}\right\rangle }{\left\Vert \mathrm{\mathbf{b}}_{i}\right\Vert ^{2}}\mathrm{\mathbf{b}}_{i}\label{eq:ProjectionResidual}
\end{align}
must therefore be orthogonal to each basis vector, and must lie in
$\mathcal{{S}}_{2}^{\bot}$, the \emph{orthogonal complement} to $\mathcal{{S}}_{2}$\footnote{The orthogonal complement of a subspace is a vector space comprised
of all vectors that are orthogonal to the subspace.}. Seen in this light, projecting a vector from $\mathcal{{S}}_{XY}$
to $\mathcal{{S}}_{Y}$ is equivalent to zeroing out its first component
and preserving its second; the residual of this projection lies entirely
in $\mathcal{{S}}_{X}$, which is the orthogonal complement to $\mathcal{{S}}_{Y}$. 

\section{Matrices and Matrix Arithmetic}

\subsection{Some Basic Definitions}

We can stack $m$ vectors\footnote{We continue to require that all components of these vectors be real
numbers, though much of this section can be generalized to complex
numbers and finite fields.} with $n$ elements each side by side to form an $n\times m$ \emph{matrix}
as follows\footnote{We follow the notation in \citet{Gentle2017}, and use $n$ to denote
the number of rows and $m$ to denote the number of columns. Other
authors exchange $n$ and $m$, but still put the number of rows before
the number of columns.}:

\begin{align}
\mathrm{\mathbf{A}} & =\left[\begin{array}{ccc}
\mathrm{\mathbf{a}}_{1} & \cdots & \mathrm{\mathbf{a}}_{m}\end{array}\right]\nonumber \\
 & =\left[\begin{array}{ccc}
a_{11} & \cdots & a_{1m}\\[0.5em]
\vdots & \vdots & \vdots\\
a_{n1} & \cdots & a_{nm}
\end{array}\right].\label{eq:Matrix}
\end{align}
If $n\neq m$, the matrix is said to be \emph{rectangular}, and if
$n=m$, it is said to be \emph{square}. Rectangular matrices often
arise when we model the returns of many securities using a few cross--sectional
or time--series factors, while square matrices arise when representing
the interaction between all possible pairs of securities or factors.
A covariance matrix, for example, is square.

The transpose of a matrix $\mathrm{\mathbf{A}}$, denoted $\mathrm{\mathbf{A}}^{\prime}$,
is obtained by exchanging its rows and columns, i.e.

\begin{equation}
\mathrm{\mathbf{A}}^{\prime}=\left[\begin{array}{ccc}
a_{11} & \cdots & a_{n1}\\
\vdots & \vdots & \vdots\\
a_{1m} & \cdots & a_{nm}
\end{array}\right].\label{eq:Atranspose}
\end{equation}

Transposition transforms an $n\times m$ matrix into a $m\times n$
matrix.

\subsection{The Rank of a Matrix}

The \emph{column space} of a matrix is the vector space spanned by
its columns, while its \emph{row space} is the vector space spanned
by its rows. The row space of a matrix is the column space of its
transpose. The \emph{null space }of a matrix\emph{, $\mathcal{{N}}(\mathrm{\mathbf{A}})$,}
as well as its \emph{left null space,} \emph{$\mathcal{{N}}(\mathrm{\mathbf{A}^{\prime}})$,}
are vector spaces that are orthogonal to its column and row spaces
respectively. They are defined by the fact that vectors in the null
space and left null space are mapped to $\mathbf{0}$ by $\mathrm{\mathbf{A}}$
and $\mathrm{\mathbf{A}}^{\prime}$ respectively, i.e.

\begin{subequations}
\begin{align}
\mathcal{{N}}\left(\mathrm{\mathbf{A}}\right) & =\left\{ \mathrm{\mathbf{x}}\,|\,\mathrm{\mathbf{A}}\mathrm{\mathbf{x}}=\mathrm{\mathbf{0}}\right\} \label{eq:NullSpace}\\
\mathcal{{N}}\left(\mathrm{\mathbf{A}}^{\prime}\right) & =\left\{ \mathrm{\mathbf{y}}\,|\,\mathrm{\mathbf{A}}^{\prime}\mathrm{\mathbf{y}}=\mathrm{\mathbf{0}}\right\} ,\label{eq:LeftNullSpace}
\end{align}
\end{subequations}

where the product of a matrix and a vector is defined in \eqref{eq:MatrixVectorProduct}.

Closely related to the dimension of a vector space (the maximum number
of linearly independent vectors contained in it) are the \emph{row
rank} of a matrix, or the maximum number of linearly independent \emph{rows}
contained in it, and its \emph{column rank}, or the maximum number
of linearly independent \emph{columns} contained in it. Although one
might think that the row and column ranks of a matrix will in general
be different, it is surprisingly true that they are always equal.
A proof of this non--intuitive result can be found in \citet{Gentle2017}.
It follows that 

\begin{equation}
\mathrm{rank}\left(\mathrm{\mathbf{A}}_{n\times m}\right)\leqslant\min\left(n,\,m\right).\label{eq:MatrixRank}
\end{equation}

Furthermore

\begin{subequations}
\begin{align}
\mathrm{rank}\left(\mathrm{\mathbf{A}}_{n\times m}\right)\,+\,\mathrm{rank}\left(\mathcal{{N}}\left(\mathrm{\mathbf{A}}_{n\times m}\right)\right) & =m\label{eq:SumOfRanksA}\\
\mathrm{rank}\left(\mathrm{\mathbf{A}^{\prime}}_{m\times n}\right)\,+\,\mathrm{rank}\left(\mathcal{{N}}\left(\mathrm{\mathbf{A}}_{m\times n}^{\prime}\right)\right) & =n,\label{eq:SumOfRanksB}
\end{align}
\end{subequations}

and 

\begin{subequations}
\begin{align}
\mathrm{rank}\left(\mathrm{\mathbf{A}}\,+\,\mathbf{B}\right) & \leqslant\mathrm{rank}\left(\mathrm{\mathbf{A}}\right)\,+\,\mathrm{rank}\left(\mathbf{B}\right)\label{eq:RankOfSums}\\
\mathrm{rank}\left(\mathrm{\mathbf{A}}\mathbf{B}\right) & \leqslant\min\left(\mathrm{rank}\left(\mathrm{\mathbf{A}}\right),\,\mathrm{rank}\left(\mathbf{B}\right)\right),\label{eq:RankOfProduct}
\end{align}
\end{subequations}

where matrix addition and multiplication are defined in \eqref{eq:MatrixAddition}
and \eqref{eq:MatrixMultiplication} respectively.


\subsection{Matrix Addition and Multiplication}

We can add two matrices elementwise if they have the same numbers
of rows and columns

\begin{equation}
\mathrm{\mathbf{A}}\,+\,\mathrm{\mathbf{B}}=\left[\begin{array}{ccc}
a_{11}\,+\,b_{11} & \ \cdots\  & a_{1m}\,+\,b_{1m}\\
\vdots & \vdots & \vdots\\
a_{n1}\,+\,b_{n1} & \ \cdots\  & a_{nm}\,+\,b_{nm}
\end{array}\right],\label{eq:MatrixAddition}
\end{equation}
and multiply a matrix by a scalar, just as we did with vectors:

\begin{equation}
c\cdotp\mathrm{\mathbf{A}}=\left[\begin{array}{ccc}
c\thinspace a_{11} & \ \cdots\  & c\thinspace a_{1m}\\
\vdots & \vdots & \vdots\\
c\thinspace a_{n1} & \ \cdots\  & c\thinspace a_{nm}
\end{array}\right].\label{eq:MatrixScaling}
\end{equation}

A matrix may also be multiplied by a vector, to give rise to another
vector. The product $\mathrm{\mathbf{Ax}}$ is a \emph{linear transformation}
of the vector $\mathrm{\mathbf{x}}$, and can also be viewed as a
linear combination of the columns of the matrix, with the weights
being given by the vector $\mathrm{\mathbf{x}}$. The number of columns
in $\mathrm{\mathbf{A}}$ must equal the number of elements in $\mathrm{\mathbf{x}}$
if the product is to be valid:

\begin{align}
\mathrm{\mathbf{Ax}} & =\left[\begin{array}{ccc}
\mathrm{\mathbf{a}}_{1} & \cdots & \mathrm{\mathbf{a}}_{m}\end{array}\right]\left[\begin{array}{c}
x_{1}\\
\vdots\\
x_{m}
\end{array}\right]\nonumber \\
 & =\left[\begin{array}{c}
x_{1}\,\mathrm{\mathbf{a}}_{1}\,+\,\cdots\,+\,x_{m}\,\mathrm{\mathbf{a}}_{m}\end{array}\right]\nonumber \\
 & =\left[\begin{array}{c}
a_{11}\thinspace x_{1}+\cdots+a_{1m}\thinspace x_{m}\\
\vdots\\
a_{n1}\thinspace x_{1}+\cdots+a_{nm}\thinspace x_{m}
\end{array}\right].\label{eq:MatrixVectorProduct}
\end{align}

As a matrix is just a stacked set of vectors, \eqref{eq:MatrixVectorProduct}
generalizes naturally to allow for the multiplication of two matrices
that are conformable -- the number of columns in the first matrix
must equal the number of rows in the second:

\begin{align}
\mathrm{\mathbf{AB}} & =\mathrm{\mathbf{A}}\left[\begin{array}{c}
\mathrm{\mathbf{b}}_{1}\cdots\mathrm{\mathbf{b}}_{p}\end{array}\right]\nonumber \\[0.5em]
 & =\left[\begin{array}{c}
\mathrm{\mathbf{A}}\mathrm{\mathbf{b}}_{1}\cdots\mathrm{\mathrm{\mathbf{A}}\mathbf{b}}_{p}\end{array}\right]\nonumber \\[0.5em]
 & =\left[\begin{array}{ccc}
a_{11}\thinspace b_{11}\,+\,\cdots\,+\,a_{1m}\thinspace b_{m1} & \ \cdots\  & a_{11}\thinspace b_{1p}\,+\,\cdots\,+\,a_{1n}\thinspace b_{mp}\\
\vdots & \vdots & \vdots\\
a_{n1}\thinspace b_{11}\,+\,\cdots\,+\,a_{nm}\thinspace b_{n1} & \ \cdots\  & a_{n1}\thinspace b_{1p}\,+\,\cdots\,+\,a_{nm}\thinspace b_{mp}
\end{array}\right].\label{eq:MatrixMultiplication}
\end{align}

If $\mathrm{\mathbf{A}}$ is $n\,\times\,m$, and $\mathrm{\mathbf{B}}$
is $m\,\times\,p$, the product $\mathrm{\mathbf{A\,B}}$ is $n\,\times\,p$.
Additionally

\begin{equation}
\left(\mathrm{\mathbf{A\,B}}\right)^{\prime}=\mathrm{\mathbf{B}}^{\prime}\,\mathrm{\mathbf{A}}^{\prime}.\label{eq:ABtranspose}
\end{equation}

We note in passing that $\mathrm{\mathbf{AB}}$ can also be viewed
as a linear combination of the rows of $\mathrm{\mathbf{B}}$, with
the weights being given by the rows of $\mathrm{\mathbf{A}}.$ In
general

\begin{equation}
\mathrm{\mathbf{A\,B}}\ne\mathrm{\mathbf{B\,A}.}\label{eq:ABneBA}
\end{equation}

Matrices prove particularly useful when solving sets of linear equations.
If we set 

\begin{equation}
\mathrm{\mathbf{A\,x}}=\mathrm{\mathbf{b},}\label{eq:AxEQb}
\end{equation}

where $\mathbf{b}$ is a column vector, we have a compact representation
of a collection of linear equations that can be interpreted in one
of two ways:
\begin{enumerate}
\item Row perspective: Each row of \eqref{eq:AxEQb} corresponds to a \emph{hyperplane}\footnote{In $n$--dimensional Euclidean space, a hyperplane is a subspace
whose dimension is one less than the dimension of the space. In 2
and 3 dimensions, hyperplanes are lines and planes respectively.}, and solving \eqref{eq:AxEQb} entails finding a vector $\mathrm{\mathbf{x}}$
that lies at the intersection of these $n$ hyperplanes
\item Column perspective: Each column of $\mathrm{\mathbf{A}}$ is a vector,
and solving \eqref{eq:AxEQb} entails finding a vector $\mathrm{\mathbf{x}}$
that makes the two sides of \eqref{eq:AxEQb} equal.
\end{enumerate}
Sometimes, \eqref{eq:AxEQb} possesses a unique solution, while at
other times, it can have an infinite number of solutions or no solution
at all. The existence of a unique solution is governed by whether
or not $\mathrm{\mathbf{A}}$ is invertible, a topic we will return
to shortly.

\subsection{Some Special Matrices}

Some matrices have a special structure that greatly simplifies operations
on them\footnote{A comprehensive list of matrices with special structures and their
names can be found at \url{https://en.wikipedia.org/wiki/List_of_named_matrices}.}. 

If a matrix is square, and if $a_{ij}=a_{ji}$ for all $i,\,j$, it
is said to be \emph{symmetric}. If, in addition, $a_{ij}=0\ \forall i\neq j$,
it is said to be \emph{diagonal}.The transpose of a symmetric matrix
is equal to itself. Clearly, all diagonal matrices are symmetric,
but the converse is not true.

The \emph{Zero matrix} $\mathrm{\mathbf{0}}_{n}$ is an $n\times n$
matrix with 0's everywhere, i.e.

\begin{equation}
\mathrm{\mathbf{0}}_{n}=\left[\begin{array}{ccc}
0 & \cdots & 0\\
\vdots & \ddots & \vdots\\
0 & \cdots & 0
\end{array}\right].\label{eq:ZeroMatrix}
\end{equation}

The \emph{Identity matrix} $\mathrm{\mathbf{I}}_{n}$ is an $n\times n$
diagonal matrix with 1's along the diagonal and 0's elsewhere, i.e.

\begin{equation}
\mathrm{\mathbf{I}}_{n}=\left[\begin{array}{ccc}
1 & \cdots & 0\\
\vdots & \ddots & \vdots\\
0 & \cdots & 1
\end{array}\right].\label{eq:IdentityMatrix}
\end{equation}

$\mathrm{\mathbf{J}}_{n}$, the $n\times n$ matrix with all 1's can
be constructed by mutliplying a vector of $n$ 1's by its transpose

\begin{align}
\mathrm{\mathbf{J}}_{n} & =\mathrm{\mathbf{1}}_{n}\,\mathrm{\mathbf{1}}_{n}^{\prime}\nonumber \\
 & =\left[\begin{array}{c}
1\\
\vdots\\
1
\end{array}\right]\left[\begin{array}{c}
1\\
\vdots\\
1
\end{array}\right]^{\prime}\nonumber \\
 & =\left[\begin{array}{ccc}
1 & \cdots & 1\\
\vdots & \ddots & \vdots\\
1 & \cdots & 1
\end{array}\right].\label{eq:AllOneMatrix}
\end{align}

An $n\times m$ matrix is said to be \emph{Orthogonal} if

\begin{equation}
\mathrm{\mathbf{A}}^{\prime}\,\mathrm{\mathbf{A}}=\mathrm{\mathbf{I}}_{m},\label{eq:OrthogonalMatrix}
\end{equation}

and it is useful to think of such a matrix as a stacked collection
of vectors that lie at right angles to each other in Euclidean space.
The product $\mathrm{\mathbf{A}}^{\prime}\mathrm{\mathbf{A}}$ is
a symmetric matrix that is known as the \emph{Gramian}, and whose
rank is the same as the rank of $\mathrm{\mathbf{A}}$.

A martrix $\mathrm{\mathbf{A}}$ can be raised to an integer power
$k$ by multiplying together $k$ copies of itself. For example

\begin{equation}
\mathrm{\mathbf{A}}^{2}=\mathrm{\mathbf{A}}\,\mathrm{\mathbf{A}},\label{eq:ASquared}
\end{equation}

and is said to be \emph{nilpotent of index k} if

\begin{subequations}
\begin{align}
\mathrm{\mathbf{A}}^{k} & =\mathbf{0}_{n}\label{eq:NilpotentMatrixA}\\
\mathrm{\mathbf{A}}^{j} & \ne\mathbf{0}_{n}\ \mathrm{for}\ 1\leqslant j<k.\label{eq:NilpotentMatrixB}
\end{align}
\end{subequations}


\subsection{Toeplitz Matrices}

A \emph{Toeplitz} matrix is a square matrix with constant offset diagonals,
i.e.

\begin{equation}
a_{i,\,j}=a_{i\,+\,1,\,j\,+\,1},\ 1\leqslant i,\,j\leqslant n-1\label{eq:ToeplitzMatrix1}
\end{equation}

e.g.,

\begin{equation}
\mathrm{\mathbf{A}}=\left[\begin{array}{ccccc}
1 & 2 & 3 & 4 & 5\\
6 & 1 & 2 & 3 & 4\\
7 & 6 & 1 & 2 & 3\\
8 & 7 & 6 & 1 & 2\\
9 & 8 & 7 & 6 & 1
\end{array}\right].\label{eq:ToeplitzMatrix2}
\end{equation}

In the special case when each row is a cyclic shift of the prior row,
the matrix is referred to as a \emph{circulant} matrix, e.g.

\begin{equation}
\mathrm{\mathbf{A}}=\left[\begin{array}{ccccc}
1 & 2 & 3 & 4 & 5\\
5 & 1 & 2 & 3 & 4\\
4 & 5 & 1 & 2 & 3\\
3 & 4 & 5 & 1 & 2\\
2 & 3 & 4 & 5 & 1
\end{array}\right].\label{eq:CirculantMatrix}
\end{equation}

Every circulant matrix is a Toeplitz matrix, but the converse is not
true.

\subsection{Block Matrices}

A matrix can often usefully be partitioned into a set of adjacent
\emph{blocks, }with the boundaries between blocks being identified
by vertical and horizontal lines. The circulant matrix displayed in
\eqref{eq:CirculantMatrix}, for example, can be partitioned into
four blocks in a variety of ways as can be seen in \eqref{eq:PartitionedMatrix1}--\eqref{eq:PartitionedMatrix4}.
Observe that the various submatrices are not necessarily square --
in particular, \emph{none} of the submatrices in \eqref{eq:PartitionedMatrix4}
are square. 

Even though a matrix can have multiple blocks, we illustrate block
matrix operations using matrices that are partitioned into four adjacent
blocks. The results presented here generalize naturally to larger
numbers of blocks, and we refer the interested reader to the comprehensive
treatment in \citet{Gentle2017}.

\begin{subequations}
\begin{align}
\mathrm{\mathbf{A}} & =\left[\begin{array}{ccccc}
1 & 2 & 3 & 4 & 5\\
5 & 1 & 2 & 3 & 4\\
4 & 5 & 1 & 2 & 3\\
3 & 4 & 5 & 1 & 2\\
2 & 3 & 4 & 5 & 1
\end{array}\right]\label{eq:PartitionedMatrix1}\\[0.5em]
 & =\left[\begin{array}{c|c}
\mathrm{\mathbf{A}}_{11}\, & \,\mathrm{\mathbf{A}}_{12}\\
\hline \mathrm{\mathbf{A}}_{21}\, & \,\mathrm{\mathbf{A}}_{22}
\end{array}\right]\label{eq:PartitionedMatrix2}\\[0.5em]
 & =\left[\begin{array}{cc|ccc}
1 & 2 & 3 & 4 & 5\\
5 & 1 & 2 & 3 & 4\\
\hline 4 & 5 & 1 & 2 & 3\\
3 & 4 & 5 & 1 & 2\\
2 & 3 & 4 & 5 & 1
\end{array}\right]\label{eq:PartitionedMatrix3}\\[0.5em]
 & =\left[\begin{array}{ccc|cc}
1 & 2 & 3 & 4 & 5\\
\hline 5 & 1 & 2 & 3 & 4\\
4 & 5 & 1 & 2 & 3\\
3 & 4 & 5 & 1 & 2\\
2 & 3 & 4 & 5 & 1
\end{array}\right]\label{eq:PartitionedMatrix4}
\end{align}
\end{subequations}

It is useful to think of the blocks along the main diagonal as representing
the influence of a tightly coupled set of items on each other (e.g.
stocks within a sector, with one block per sector), while blocks off
the main diagonal represent interactions between different sets of
items (i.e. between sectors in an index). Sometimes, the off--diagonal
blocks have entries that are much smaller in magnitude than the blocks
on the diagonal, and may be set to 0 to simplify the matrix. A matrix
whose off--diagonal blocks are all zero is called a \emph{block diagonal}
matrix.

The transpose of a block matrix is obtained by block transposition
of the transposes of each of the blocks:

\begin{align}
\mathrm{\mathbf{A}}^{\prime} & =\left[\begin{array}{c|c}
\mathrm{\mathbf{A}}_{11} & \mathrm{\mathbf{A}}_{12}\\
\hline \mathrm{\mathbf{A}}_{21} & \mathrm{\mathbf{A}}_{22}
\end{array}\right]^{\prime}\nonumber \\[0.5em]
 & =\left[\begin{array}{c|c}
\mathrm{\mathbf{A}}_{11}^{\prime} & \mathrm{\mathbf{A}}_{21}^{\prime}\\
\hline \mathrm{\mathbf{A}}_{12}^{\prime} & \mathrm{\mathbf{A}}_{22}^{\prime}
\end{array}\right].\label{eq:BlockTranspose}
\end{align}

Just as matrices allow addition and multiplication if their dimensions
are conformable, so do block matrices allow block addition and block
multiplication provided that the dimensions of the various blocks
are conformable, e.g. 

\begin{align}
\mathrm{\mathbf{A+B}} & =\left[\begin{array}{c|c}
\mathrm{\mathbf{A}}_{11}\, & \,\mathrm{\mathbf{A}}_{12}\\
\hline \mathrm{\mathbf{A}}_{21}\, & \,\mathrm{\mathbf{A}}_{22}
\end{array}\right]+\left[\begin{array}{c|c}
\mathrm{\mathbf{B}}_{11}\, & \,\mathrm{\mathbf{B}}_{12}\\
\hline \mathrm{\mathbf{B}}_{21}\, & \,\mathrm{\mathbf{B}}_{22}
\end{array}\right]\nonumber \\[0.5em]
 & =\left[\begin{array}{c|c}
\mathrm{\mathbf{A}}_{11}\,+\,\mathrm{\mathbf{B}}_{11}\, & \,\mathrm{\mathbf{A}}_{12}\,+\,\mathrm{\mathbf{B}}_{12}\\
\hline \mathrm{\mathbf{A}}_{21}\,+\,\mathrm{\mathbf{B}}_{21}\, & \,\mathbf{A}_{22}\,+\,\mathrm{\mathbf{B}}_{22}
\end{array}\right],\label{eq:BlockMatrixAddition}
\end{align}

and

\begin{align}
\mathrm{\mathbf{A\,B}} & =\left[\begin{array}{c|c}
\mathrm{\mathbf{A}}_{11}\, & \,\mathrm{\mathbf{A}}_{12}\\
\hline \mathrm{\mathbf{A}}_{21}\, & \,\mathrm{\mathbf{A}}_{22}
\end{array}\right]\left[\begin{array}{c|c}
\mathrm{\mathbf{B}}_{11}\, & \,\mathrm{\mathbf{B}}_{12}\\
\hline \mathrm{\mathbf{B}}_{21}\, & \,\mathrm{\mathbf{B}}_{22}
\end{array}\right]\nonumber \\[0.5em]
 & =\left[\begin{array}{c|c}
\mathrm{\mathbf{A}}_{11}\mathrm{\mathbf{B}}_{11}\,+\,\mathrm{\mathbf{A}}_{12}\mathrm{\mathbf{B}}_{21}\, & \,\mathrm{\mathbf{A}}_{11}\mathrm{\mathbf{B}}_{12}\,+\,\mathrm{\mathbf{A}}_{12}\mathrm{\mathbf{B}}_{22}\\
\hline \mathrm{\mathbf{A}}_{21}\mathrm{\mathbf{B}}_{11}\,+\,\mathrm{\mathbf{A}}_{22}\mathrm{\mathbf{B}}_{21}\, & \,\mathrm{\mathbf{A}}_{21}\mathrm{\mathbf{B}}_{12}\,+\,\mathrm{\mathbf{A}}_{22}\mathrm{\mathbf{B}}_{22}
\end{array}\right].\label{eq:BlockMatrixMultiplication}
\end{align}

The partitioned matrix in \eqref{eq:PartitionedMatrix3} can be multiplied
by itself in accordance with \eqref{eq:BlockMatrixMultiplication},
but the partition in \eqref{eq:PartitionedMatrix4} cannot be block
multiplied by itself or by the partition in \eqref{eq:PartitionedMatrix3}.


\subsection{Echelon Form and LU Factorization}

Computing the rank of a matrix has much in common with solving the
set of linear equations specified in \eqref{eq:AxEQb}. Using Gaussian
elimination\footnote{Gaussian elimination attempts to triangularize a matrix by iteratively
subtracting a multiple of each row, starting from the first row and
working downwards, from all the rows below so as to cancel their leading
entry. The process transforms a matrix into row echelon form, and
ideally into an upper triangular matrix.}, one can transform $\mathrm{\mathbf{A}}$ into a matrix $\mathbf{E}$
in \emph{row echelon form} with the following properties
\begin{enumerate}
\item All non-zero rows of $\mathbf{E}$ lie above all rows with only zeros
\item The first non--zero entry in each row lies strictly to the right
of the first non--zero entry in the row above it, and
\item All entries below the first non--zero entry in each row are 0.
\end{enumerate}
The rank of $\mathbf{A}$ is the number of non--zero rows in $\mathbf{E}$.
It is also the dimension of its column and row spaces. The sum of
the dimensions of the column space and the null space must equal the
number of columns, while the sum of the dimensions of the row space
and the left null space must equal the number of rows.

If $\mathrm{\mathbf{A}}$ is square and of full rank, it can be represented
as the product of lower and upper triangular matrices $\mathrm{\mathbf{L}}$
and $\mathrm{\mathbf{U}}$, i.e.

\begin{equation}
\mathrm{\mathbf{A}}=\mathrm{\mathbf{L}}\,\mathrm{\mathbf{U},}\label{eq:AequalsLU}
\end{equation}

where $\mathrm{\mathbf{L}}$ has only 1's on the diagonal, permitting
a rapid solution to \eqref{eq:AxEQb}. A particular useful form of
the LU factorization is the LDU factorization

\begin{equation}
\mathrm{\mathbf{A}}=\mathrm{\mathbf{L}}\,\mathrm{\mathbf{D}}\,\mathbf{U},\label{eq:AequalsLDU}
\end{equation}

where $\mathrm{\mathbf{D}}$ is a diagonal matrix, and both $\mathrm{\mathbf{L}}$
and $\mathrm{\mathbf{U}}$ have only 1's on the diagonal. 

\subsection{The Determinant of a Matrix}

The determinant of a square matrix $\mathrm{\mathbf{A}}$, denoted
$det\left(\mathrm{\mathbf{A}}\right)$, is the signed volume of a
parallelepiped whose vertices are the $2^{n}$ possible sums of the
rows of $\mathrm{\mathbf{A}}$\footnote{The origin, which corresponds to the sum of none of the rows, is one
of the vertices of the parallelepiped.}. Unlike the volume of a physical object, which is always positive,
the determinant of a matrix has both sign and magnitude. In this regard,
it has something in common with vectors. If $det\left(\mathrm{\mathbf{A}}\right)=0$,
$\mathrm{\mathbf{A}}$ is said to be \emph{singular}. 

It is particularly easy to compute the determinant of a diagonal matrix
$\mathrm{\mathbf{D}}$, as its $2^{n}$ possible row sums represent
an $n-$dimensional hypercube with orthogonal sides. Its volume (and
determinant) are therefore given by

\begin{equation}
det\left(\mathrm{\mathbf{D}}\right)=\prod_{i\,=\,1}^{n}d_{ii}.\label{eq:DiagonalDeterminant}
\end{equation}

The row sums of the $n\times n$ identity matrix $\mathrm{\mathbf{I}}_{n}$
represent a unit hypercube, and its determinant must therefore be
1. 

More generally, the determinant of a block diagonal matrix is the
product of the determinants of each of its diagonal blocks, i.e.
\begin{equation}
det\,\left[\begin{array}{cccc}
\mathrm{\mathbf{D}}_{11} & \mathrm{\mathbf{0}} & \cdots & \mathrm{\mathbf{0}}\\
\mathrm{\mathbf{0}} & \ddots &  & \mathrm{\mathbf{0}}\\
\vdots &  &  & \vdots\\
\mathrm{\mathbf{0}} & \mathrm{\mathbf{0}} & \cdots & \mathrm{\mathbf{D}}_{nn}
\end{array}\right]=\prod_{i\,=\,1}^{n}det\left(\mathrm{\mathbf{D}}_{ii}\right).\label{eq:BlockDiagonalDeterminant}
\end{equation}

A $2\times2$ matrix vividly illustrates many properties of determinants,
though it is not always possible to project the insights it affords
to understand larger matrices. It is easy to show that the area of
a quadrilateral with vertices $(0,\,0)$, $(a,\,b)$, $(c,\,d)$ and
$(a+c,\,b+d)$ is $ad-bc$, and it follows that

\begin{equation}
det\mathrm{\,\left[\begin{array}{cc}
\mathit{a}\, & \,\mathit{b}\\
\mathit{c}\, & \,\mathit{d}
\end{array}\right]}=ad\,-\,bc.\label{eq:Det2by2}
\end{equation}

If any two rows of $\mathrm{\mathbf{A}}$ are interchanged, the magnitude
of its determinant remains unchanged, but its sign reverses -- indicating
that we are visiting the vertices of the hypercube in a different
sequence, e.g.

\begin{align}
det\mathrm{\,\left[\begin{array}{cc}
\mathit{c}\, & \,\mathit{d}\\
\mathit{a}\, & \text{\,}\mathit{b}
\end{array}\right]} & =bc\,-\,ad\nonumber \\[0.5em]
 & =-\,det\mathrm{\,\left[\begin{array}{cc}
\mathit{a}\, & \,\mathit{b}\\
\mathit{c}\, & \,\mathit{d}
\end{array}\right]}.\label{eq:RowExchange2by2}
\end{align}

Furthermore, if any two rows of $\mathrm{\mathbf{A}}$ are identical,
its determinant must be 0. This follows from the observation that
exchanging two rows must reverse the sign of the determinant. But
if the rows are identical, the determinant must remain the same, and
this can only happen if it is 0.

Multiplying any single row of $\mathrm{\mathbf{A}}$ by a constant
$c$ scales its determinant by $c$ as well -- the volume of a hypercube
is proportional to the length of each of its sides. It follows that
if the entire matrix is multiplied by $c$, its determinant is scaled
by $c^{n}$, i.e.

\begin{equation}
det\left(c\mathrm{\mathbf{A}}\right)=c^{n}det\left(\mathrm{\mathbf{A}}\right).\label{eq:ScaledMatixDeterminant}
\end{equation}

If any row of $\mathrm{\mathbf{A}}$ is 0, so is its determinant,
as this is equivalent to multiplying that row by 0.

Matrix transposition does not impact the value of the determinant,
i.e.

\begin{equation}
det\left(\mathrm{\mathbf{A}^{\prime}}\right)=det\left(\mathrm{\mathbf{A}}\right),\label{eq:DeterminantofTranspose}
\end{equation}

and by way of example, it is easy to check that the area of a quadrilateral
with vertices $(0,\,0)$, $(a,\,d)$, $(b,\,c)$ and $(a+b,\,c+d)$
is $ad-bc$. This is \emph{not} the same quadrilateral as that associated
with \eqref{eq:Det2by2}, but their areas are identical.

Additionally, adding a multiple of any row to any other row (or any
column to any other column) does not change the value of the determinant,
e.g.

\begin{align}
det\,\left[\begin{array}{cc}
\mathit{a} & \mathit{b}\\
\mathit{c}\,+\,na\, & \,\mathit{d\,+\,nb}
\end{array}\right] & =a\left(d\,+\,nb\right)\,-\,b\left(c\,+\,na\right)\nonumber \\[0.5em]
 & =ad\,-\,bc\nonumber \\[0.5em]
 & =det\,\left[\begin{array}{cc}
\mathit{a}\, & \,\mathit{b}\,+\,ma\\
\mathit{c}\, & \,\mathit{d\,+\,mc}
\end{array}\right]\nonumber \\[0.5em]
 & =det\mathrm{\,\left[\begin{array}{cc}
\mathit{a}\, & \,\mathit{b}\\
\mathit{c}\, & \,\mathit{d}
\end{array}\right]}.\label{eq:RowAddition2by2}
\end{align}

In particular, if $a\ne0$ and $n=-\nicefrac{c}{a}$ in \eqref{eq:RowAddition2by2},
the matrix becomes upper triangular, and its determinant is just the
product of its diagonal entries. 

\begin{equation}
det\mathrm{\,\left[\begin{array}{cl}
\mathit{a}\, & \,\mathit{b}\\
0\, & \,\mathit{\frac{ad\,-\,bc}{a}}
\end{array}\right]}=ad\,-\,bc.\label{eq:Det2by2Triangular}
\end{equation}

More generally, 

\begin{align}
det\,\left[\begin{array}{cccc}
\mathit{t_{11}} & t_{12} & \cdots & t_{1n}\\
0 & t_{22} &  & t_{2n}\\
\vdots & \vdots & \ddots & \vdots\\
0 & 0 & \cdots & t_{nn}
\end{array}\right] & =det\,\left[\begin{array}{cccc}
\mathit{t_{11}} & 0 & \cdots & 0\\
t_{21} & t_{22} &  & 0\\
\vdots & \vdots & \ddots & \vdots\\
t_{n1} & t_{n2} & \cdots & t_{nn}
\end{array}\right]\nonumber \\[0.5em]
 & =\prod_{i\,=\,1}^{n}t_{ii}.\label{eq:TriangularMatrixDerminant}
\end{align}

This follows from the observation that elementary row operations can
be used to to combine rows and diagonalize a triangular matrix without
changing its determinant. In a like manner, the determinant of a block
triangular matrix is the product of the determinants of its diagonal
blocks:

\begin{align}
det\,\left[\begin{array}{cccc}
\mathrm{\mathbf{T}}_{11} & \mathrm{\mathbf{T}}_{12} & \cdots & \mathrm{\mathbf{T}}_{1n}\\
\mathrm{\mathbf{0}} & \mathrm{\mathbf{T}}_{22} &  & \mathrm{\mathbf{T}}_{2n}\\
\vdots &  & \ddots & \vdots\\
\mathbf{0} & \mathbf{0} & \cdots & \mathrm{\mathbf{T}}_{nn}
\end{array}\right] & =det\,\left[\begin{array}{cccc}
\mathrm{\mathbf{T}}_{11} & \mathrm{\mathbf{0}} & \cdots & \mathrm{\mathbf{0}}\\
\mathrm{\mathbf{T}}_{21} & \mathrm{\mathbf{T}}_{22} &  & \mathrm{\mathbf{0}}\\
\vdots &  & \ddots & \vdots\\
\mathrm{\mathbf{T}}_{n1} & \mathrm{\mathbf{T}}_{n2} & \cdots & \mathrm{\mathbf{T}}_{nn}
\end{array}\right]\nonumber \\[0.5em]
 & =\prod_{i\,=\,1}^{n}det\left(\mathrm{\mathbf{T}}_{ii}\right).\label{eq:BlockTriangularDeterminant}
\end{align}

Determinants are multiplicative:

\begin{equation}
det\left(\mathbf{A}\mathbf{B}\right)=det\left(\mathrm{\mathbf{A}}\right)\,det\left(\mathrm{\mathbf{B}}\right),\label{eq:DeterminantOfProducts}
\end{equation}

and, as a consequence, the LDU factorization of a matrix, which is
of great practical importance, preserves the value of its determinant.
The determinants of all three matrices are given by the product of
their diagonal entries, and the determinants of both $\mathrm{\mathbf{L}}$
and $\mathrm{\mathbf{U}}$ are 1, so that the determinant of $\mathrm{\mathbf{D}}$
must be identical to that of $\mathrm{\mathbf{A}}$.

\subsection{The Inverse of a Matrix}

If an $n\times n$ square matrix $\mathbf{A}$ is of full rank, its
determinant is non--zero, and it is said to \emph{invertible}; there
exists a unique $n\times n$ square matrix, denoted $\mathbf{A}^{-1}$,
such that

\begin{align}
\mathbf{A}^{-1}\,\mathbf{A} & =\mathbf{A}\,\mathbf{A}^{-1}\nonumber \\
 & =\mathbf{I}_{n}.\label{eq:AinverseAeqI}
\end{align}

Matrix inversion can be thought of as the generalization of the reciprocal
to matrices, and this interpretation is clearly visible when inverting
a diagonal matrix . If

\begin{align}
\mathrm{\mathbf{D}} & =\left[\begin{array}{cccc}
d_{11} & 0 & \cdots & 0\\
0 & \ddots &  & \vdots\\
\vdots &  & \ddots & 0\\
0 & \cdots & 0 & d_{nn}
\end{array}\right]\nonumber \\[0.5em]
\mathrm{\mathbf{D}}^{-1} & =\left[\begin{array}{cccc}
\frac{1}{d_{11}} & 0 & \cdots & 0\\
0 & \ddots &  & \vdots\\
\vdots &  & \ddots & 0\\
0 & \cdots & 0 & \frac{1}{d_{nn}}
\end{array}\right]\nonumber \\
 & =\frac{1}{det\left(\mathrm{\mathbf{D}}\right)}\left[\begin{array}{cccc}
\prod_{i\,\ne\,1}d_{ii} & 0 & \cdots & 0\\
0 & \ddots &  & \vdots\\
\vdots &  & \ddots & 0\\
0 & \cdots & 0 & \prod_{i\,\ne\,n}d_{ii}
\end{array}\right].\label{eq:Dinverse}
\end{align}

In general, the inverse of a matrix is given by the ratio of the transpose
of its \emph{cofactor matrix} \footnote{We use the letter $\mathrm{\mathbf{F}}$ to denote the cofactor matrix
as $\mathrm{\mathbf{C}}$ is used to denote the covariance matrix.} and its determinant

\begin{equation}
\mathrm{\mathbf{A}}^{-1}=\frac{\mathrm{\mathbf{F}}^{\prime}}{det\left(\mathrm{\mathbf{A}}\right)},\label{eq:CofactorExpansion}
\end{equation}

where $f_{ij}$ is $\left(-1\right)^{i+j}$ times the determinant
of $\mathrm{\mathbf{A}}$ with its $i^{th}$ row and $j^{th}$column
removed. The inverse of a $2\times2$ matrix, for example, is

\begin{equation}
\left[\begin{array}{cc}
a\, & \,b\\
c\, & \,d
\end{array}\right]^{-1}=\frac{1}{ad\,-\,bc}\left[\begin{array}{rr}
d\, & \,-b\\
-c\, & \,a
\end{array}\right].\label{eq:TwoByTwoInverse}
\end{equation}

The inverse of a product of two matrices exhibits the same ``reverse
order'' property we saw earlier for its transpose:

\begin{equation}
\left(\mathrm{\mathbf{A\,B}}\right)^{-1}=\mathrm{\mathbf{B}}^{-1}\,\mathrm{\mathbf{A}}^{-1},\label{eq:ABinverse}
\end{equation}

and the availability of an inverse allows a closed form solution to
\eqref{eq:AxEQb}:

\begin{equation}
\mathrm{\mathbf{x}}=\mathrm{\mathbf{A}}^{-1}\,\mathrm{\mathbf{b}}.\label{eq:xEqAinverseB}
\end{equation}

Somewhat surprisingly, and in spite of its great conceptual importance,
matrix inversion is rarely used when solving linear equations on account
of the fact that the the computation of the cofactor matrix is both
tedious and numerically unstable. Instead, Gaussian elimination and
LU factorization are used to solve \eqref{eq:AxEQb} directly.

\subsection{Inverses of Special Matrices }

There appears to be a virtually inexhaustible supply of results involving
the inverses of block matrices of various types -- the results we
present here are those most likely to prove useful in portfolio construction
and risk management and are drawn from \citet{LuShiou2002}. $2\times2$
block matrices prove particularly useful, as they are often used to
split securities into two subsets (loosely speaking, a set of securities
of interest and another set containing all other securities). Formulas
involving the inverse of a submatrix are valid only if the inverse
exists. Likewise, formulas that involve the product of matrices are
valid only if the relevant submatrices are conformable.

We start by stating results for the inverses of diagonal and anti--diagonal
block matrices\footnote{We take the liberty of naming the blocks $\mathrm{\mathbf{A}},\ \mathbf{B},\ \mathrm{\mathbf{C}}$
and $\mathrm{\mathbf{D}}$ and not $\mathrm{\mathbf{A}}_{11},\ \mathrm{\mathbf{A}}_{12},\ \mathrm{\mathbf{A}}_{21}$
and $\mathrm{\mathbf{A}}_{22}$ solely for clarity of exposition.}, and indicate their sizes via subscripts:

\begin{subequations}
\begin{align}
\left[\begin{array}{rl}
\mathbf{A}_{n\,\times\,n} & \hphantom{-}\mathbf{0}_{n\,\times\,m}\\
\mathbf{0}_{m\,\times\,n} & \boldsymbol{\hphantom{-}\mathbf{B}}_{m\,\times\,m}
\end{array}\right]^{-1} & =\left[\begin{array}{rl}
\mathbf{A}_{n\,\times\,n}^{-1} & \hphantom{-}\mathbf{0}_{n\,\times\,m}\\
\mathbf{0}_{m\,\times\,n} & \hphantom{-}\mathbf{B}_{m\,\times\,m}^{-1}
\end{array}\right]\label{eq:SpecialInverse1a}\\[0.5em]
\left[\begin{array}{lr}
\mathbf{0}_{n\,\times\,m} & \hphantom{-}\mathbf{A}_{n\,\times\,n}\\
\mathbf{B}_{m\,\times\,m} & \hphantom{-}\mathbf{0}_{m\,\times\,n}
\end{array}\right]^{-1} & =\left[\begin{array}{rl}
\mathbf{0}_{m\,\times\,n} & \hphantom{-}\mathbf{B}_{m\,\times\,m}^{-1}\\
\mathbf{A}_{n\,\times\,n}^{-1} & \hphantom{-}\mathbf{0}_{n\,\times\,m}
\end{array}\right].\label{eq:SpecialInverese1b}
\end{align}
\end{subequations}

We next provide expressions for the inverses of upper and lower triangular
matrices that are valid if $\mathrm{\mathbf{A}}$ and $\mathrm{\mathbf{C}}$
are invertible:

\begin{subequations}
\begin{align}
\left[\begin{array}{ll}
\mathrm{\mathbf{A}}_{n\,\times\,n} & \hphantom{-}\mathrm{\mathbf{B}}_{n\,\times\,m}\\
\mathbf{0}_{m\,\times\,n} & \hphantom{-}\mathrm{\mathbf{C}}_{m\,\times\,m}
\end{array}\right]^{-1} & =\left[\begin{array}{ll}
\mathbf{A}^{-1} & \mathrm{\hphantom{-}-\mathbf{A}^{-1}}\,\mathbf{B}\,\mathrm{\mathbf{C}}^{-1}\\
\mathbf{0}_{m\times n} & \hphantom{-}\mathrm{\hphantom{-}}\mathbf{C}^{-1}
\end{array}\right]\label{eq:SpecialInverse2a}\\[0.5em]
\left[\begin{array}{ll}
\mathrm{\mathbf{A}}_{n\,\times\,n} & \hphantom{-}\mathbf{0}_{n\,\times\,m}\\
\mathrm{\mathbf{B}}_{m\,\times\,n} & \hphantom{-}\mathrm{\mathbf{C}}_{m\,\times\,m}
\end{array}\right]^{-1} & =\left[\begin{array}{ll}
\hphantom{-}\mathrm{\mathbf{\mathbf{A}^{-1}}} & \hphantom{-}\mathbf{0}_{n\times m}\\
-\mathrm{\mathbf{C}}^{-1}\,\mathbf{B}\,\mathbf{A}^{-1} & \mathrm{\hphantom{-}}\mathbf{C}^{-1}
\end{array}\right].\label{eq:SpecialInverse2b}
\end{align}
\end{subequations}

In the important special case when $\mathrm{\mathbf{A}}=\mathrm{\mathbf{I}}_{n}$
and $\mathrm{\mathbf{C}}=\mathrm{\mathbf{I}}_{m}$, \eqref{eq:SpecialInverse2a}
and \eqref{eq:SpecialInverse2b} simplify to

\begin{subequations}
\begin{align}
\left[\begin{array}{ll}
\mathrm{\mathbf{I}}_{n} & \hphantom{-}\mathrm{\mathbf{B}}_{n\,\times\,m}\\
\mathbf{0}_{m\,\times\,n} & \hphantom{-}\mathrm{\mathbf{I}}_{m}
\end{array}\right]^{-1} & =\left[\begin{array}{ll}
\mathrm{\mathbf{I}}_{n} & \mathrm{\hphantom{-}-\mathbf{B}}_{n\,\times\,m}\\
\mathbf{0}_{m\,\times\,n} & \hphantom{-}\mathrm{\hphantom{-}\mathbf{I}}_{m}
\end{array}\right]\label{eq:SpecialInverse3a}\\[0.5em]
\left[\begin{array}{ll}
\mathrm{\mathbf{I}}_{n} & \hphantom{-}\mathbf{0}_{n\,\times\,m}\\
\mathrm{\mathbf{B}}_{m\,\times\,n} & \hphantom{-}\mathrm{\mathbf{I}}_{m}
\end{array}\right]^{-1} & =\left[\begin{array}{ll}
\hphantom{-}\mathrm{\mathbf{I}}_{n} & \hphantom{-}\hphantom{-}\mathbf{0}_{n\,\times\,m}\\
-\mathrm{\mathbf{B}}_{m\,\times\,n} & \hphantom{-}\mathrm{\hphantom{-}\mathbf{I}}_{m}
\end{array}\right].\label{eq:SpecialInverse3b}
\end{align}
\end{subequations}

In the general case when all 4 matrices are non-zero, the inverse
is obtained by inverting each of the matrices in a block LDU or UDL
representations using \eqref{eq:SpecialInverse1a}, \eqref{eq:SpecialInverse3a}
and \eqref{eq:SpecialInverse3b}, and finally combining their inverses
using \eqref{eq:ABinverse}:

\begin{subequations}
\begin{align}
\left[\begin{array}{cc}
\mathrm{\mathbf{A}}_{n\,\times\,n} & \mathrm{\mathbf{B}}_{n\,\times\,m}\\
\mathrm{\mathbf{C}}_{m\,\times\,n} & \mathrm{\mathbf{D}}_{m\,\times\,m}
\end{array}\right] & =\left[\begin{array}{ll}
\mathrm{\mathbf{I}}_{n} & \hphantom{-}\mathrm{\mathbf{0}}_{n\,\times\,m}\\
\mathrm{\mathbf{C}\,}\mathrm{\mathbf{A}}^{-1} & \hphantom{-}\mathrm{\mathbf{I}}_{m}
\end{array}\right]\left[\begin{array}{ll}
\mathrm{\mathbf{A}} & \hphantom{-}\mathrm{\mathbf{0}}_{n\,\times\,m}\\
\mathrm{\mathbf{0}}_{m\,\times\,n} & \mathrm{\hphantom{-}\mathbf{D}}\,-\,\mathrm{\mathbf{C}}\,\mathrm{\mathbf{A}}^{-1}\,\mathrm{\mathbf{B}}
\end{array}\right]\left[\begin{array}{ll}
\mathrm{\mathbf{I}}_{n} & \hphantom{-}\mathrm{\mathbf{A}}^{-1}\,\mathrm{\mathbf{B}}\\
\mathrm{\mathbf{0}}_{m\,\times\,n} & \hphantom{-}\mathrm{\mathbf{I}}_{m}
\end{array}\right]\label{eq:SpecialInverse4a}\\[0.5em]
 & =\left[\begin{array}{ll}
\mathrm{\mathbf{I}}_{n} & \hphantom{-}\mathrm{\mathbf{B}}\,\mathrm{\mathbf{D}}^{-1}\\
\mathrm{\mathbf{0}}_{m\,\times\,n} & \hphantom{-}\mathrm{\mathbf{I}}_{m}
\end{array}\right]\left[\begin{array}{ll}
\mathrm{\mathbf{A}\,-\,\mathrm{\mathbf{B}}\,\mathrm{\mathbf{D}}^{-1}\,\mathrm{\mathbf{C}}} & \hphantom{-}\mathrm{\mathbf{0}}_{n\,\times\,m}\\
\mathrm{\mathbf{0}}_{m\,\times\,n} & \hphantom{-}\mathbf{D}
\end{array}\right]\left[\begin{array}{ll}
\mathrm{\mathbf{I}}_{n} & \hphantom{-}\mathrm{\mathbf{0}}_{n\,\times\,m}\\
\mathrm{\mathbf{D}}^{-1}\,\mathrm{\mathbf{C}} & \hphantom{-}\mathrm{\mathbf{I}}_{m}
\end{array}\right].\label{eq:SpecialInverse4b}
\end{align}
\end{subequations}

The choice of which expansion to use will depend on whether $\mathrm{\mathbf{A}}$
or $\mathrm{\mathbf{D}}$ is invertible. Furthermore, if the matrix
is to be inverted, either $\mathrm{\mathbf{D}}-\mathrm{\mathbf{C}}\mathrm{\mathbf{A}}^{-1}\mathrm{\mathbf{B}}$
or $\mathrm{\mathbf{A}-\mathrm{\mathbf{B}}\mathrm{\mathbf{D}}^{-1}\mathrm{\mathbf{C}}}$
must additionally be invertible.

While formulae are known for the inverse of Toeplitz matrices (see
\citet{LvHuang2007}), they are cumbersome. In general, the inverse
of a Toeplitz matrix is not necessarily a Toeplitz matrix, but the
inverse of a circulant matrix is also a circulant matrix. One Toeplitz
matrix that sees extensive use in time series work and its inverse
are 

\begin{equation}
\left[\begin{array}{cccccc}
1 & \rho & \rho^{2} & \cdots &  & \rho^{n-1}\\
\rho & 1 & \rho & \cdots &  & \rho^{n-2}\\
\vdots & \vdots & \ddots & \vdots & \vdots & \vdots\\
\vdots & \vdots & \vdots & \ddots & \vdots & \vdots\\
\rho^{n-2} & \rho^{n-3} & \rho^{n-4} & \cdots & 1 & \rho\\
\rho^{n-1} & \rho^{n-2} & \rho^{n-3} & \cdots & \rho & 1
\end{array}\right]^{-1}=\left[\begin{array}{cccccc}
1 & -\rho & 0 & 0 & \cdots & 0\\
-\rho & \,1\,+\,\rho^{2}\, & -\rho & 0 & \cdots & 0\\
0 & -\rho & \,1\,+\,\rho^{2}\, & -\rho & \vdots & \vdots\\
\vdots & \vdots & -\rho & \,1\,+\,\rho^{2}\, & -\rho & 0\\
0 & 0 & \cdots & -\rho & \,1\,+\,\rho^{2}\, & -\rho\\
0 & 0 & \cdots & 0 & -\rho & 1
\end{array}\right].\label{eq:ToeplitzInverse}
\end{equation}


\subsubsection*{}

\subsection{The Sherman--Morrison--Woodbury Matrix Inversion Lemma}

The following identities see wide application in factor models, as
well as in statistics and signal processing. They are variously known
as the Sherman--Morrison--Woodbury matrix inversion lemma, the Sherman--Morrison
formula, and the Woodbury matrix identity. The fact that they are
attributed to many authors is likely a reflection of the fact that
they were discovered independently by a multitude of researchers at
about the same time. 

Broadly speaking, they provide an affirmative answer to the following
question: If, after computing the inverse of $\mathrm{\mathbf{A}}$,
we perturb it with a low rank (as opposed to a low magnitude) perturbation,
can we efficiently compute the inverse of the perturbed matrix by
perturbing $\mathrm{\mathbf{A}}^{-1}$ with a related low rank perturbation\footnote{Proofs of these, and other related, identities can be found at \url{https://en.wikipedia.org/wiki/Woodbury_matrix_identity}.}? 

We start with the special case when the perturbation has rank 1. 

Let $\mathbf{A}$ be an invertible $n\times n$ matrix, let $\mathbf{x}$
and $\mathbf{y}$ be column vectors of length $n$ (and hence of rank
1), and let $c$ be any scalar. Then:

\begin{equation}
\left(\mathbf{A}\,+\,c\,\mathbf{x}\,\mathbf{y}^{\prime}\right)^{-1}=\mathbf{A}^{-1}\,-\,\frac{c\,\mathbf{A}^{-1}\,\mathbf{x}\,\mathbf{y}^{\prime}\,\mathbf{A}^{-1}}{1\,+\,c\,\mathbf{x}^{\prime}\,\mathbf{A}^{-1}\mathbf{y}}.\label{eq:ShermanMorrison1}
\end{equation}

Equation \eqref{eq:ShermanMorrison1} simplifies even further when
$\mathrm{\mathbf{A}}=\mathrm{\mathbf{I}}_{n}$, for in this case,
$\mathbf{A}^{-1}=\mathbf{A}$, and we get

\begin{equation}
\left(\mathbf{I}_{n}\,+\,c\,\mathbf{x}\,\mathbf{y}^{\prime}\right)^{-1}=\mathbf{I}_{n}\,-\,\frac{c\mathbf{\,x}\,\mathbf{y}^{\prime}}{1\,+\,c\,\mathbf{x}^{\prime}\,\mathbf{y}}.\label{eq:ShermanMorrison2}
\end{equation}

Both \eqref{eq:ShermanMorrison1} and \eqref{eq:ShermanMorrison2}
can be evaluated in $O\left(n^{2}\right)$ arithmetic operations (direct
matrix inversion using Gaussian elimination requires $O\left(n^{3}\right)$
operations), and even more importantly, they are numerically stable
-- small changes to any of the parameters on the left hand sides
of \eqref{eq:ShermanMorrison1} and \eqref{eq:ShermanMorrison2} result
in small changes to their right hand sides. 

More generally, let $\mathbf{A}$ be $n\times n$, $\mathbf{U}$ be
$n\times k$, $\mathbf{C}$ be $k\times k$, and $\mathbf{V}$ be
$k\times n$. Then

\begin{equation}
\left(\mathbf{A}\,+\,\mathbf{U\,\mathbf{C\mathbf{\,V}}}\right)^{-1}=\mathbf{A}^{-1}\,-\,\mathbf{A}^{-1}\,\mathbf{U}\left(\mathbf{C}^{-1}\,+\,\mathbf{V}\,\mathbf{A}^{-1}\,\mathbf{U}\right)^{-1}\,\mathbf{V}\,\mathbf{A}^{-1}\label{eq:Woodbury}
\end{equation}

where it is assumed that the inverses all exist. Typically, $k\ll n$,
making the inversion of $\mathrm{\mathbf{C}}$ far less tedious than
direct recomputation of $\mathbf{A}^{-1}$.

\citet{HendersonSearle1981} and \citet{Hager1989} provide a useful
compendium of results related to \eqref{eq:ShermanMorrison1}, \eqref{eq:ShermanMorrison2},
and \eqref{eq:Woodbury}, as well as a comprehensive summary of how
they are applied in a variety of settings. One such result is

\begin{align}
\left[\begin{array}{cccc}
a & b & \cdots & b\\
b & a & \cdots & b\\
\vdots & \vdots & \ddots & \vdots\\
b & b & \cdots & a
\end{array}\right]^{-1} & =\left[\left(a\,-\,b\right)\mathrm{\mathbf{I}}_{n}\,+\,b\,\mathbf{1}_{n}\,\mathbf{1}_{n}^{\prime}\right]^{-1}\nonumber \\
 & =\frac{1}{a\,-\,b}\left[\mathrm{\mathbf{I}}_{n}\,-\,\frac{b}{a\,+\,\left(n\,-\,1\right)\,b}\,\mathbf{J}_{n}\right].\label{eq:SingleFactorShermanMorrison}
\end{align}

where $\mathbf{J}_{n}=\mathbf{1}_{n}\mathbf{1}_{n}^{\prime}$ is defined
in \eqref{eq:AllOneMatrix}, and provided that $a\,\ne\,b$ and $a\,+\,(n-1)\,b\ne0$. 

This result proves particularly useful when analyzing constant correlation
matrices with $a=1$ and $b=\rho$ for $\nicefrac{-1\,}{\,\left(n-1\right)}<\rho<1$\footnote{Outside this range, the matrix may be invertible, but it is not positive
definite.}, i.e.

\begin{align}
\left[\begin{array}{cccc}
1 & \rho & \cdots & \rho\\
\rho & 1 & \cdots & \rho\\
\vdots & \vdots & \ddots & \vdots\\
\rho & \rho & \cdots & 1
\end{array}\right]^{-1} & =\frac{1}{1\,-\,\rho}\left[\mathrm{\mathbf{I}}_{n}\,-\,\frac{\rho}{1\,+\,\left(n\,-\,1\right)\,\rho}\,\mathbf{J}_{n}\right]\nonumber \\
 & =\frac{\left(1\,+\,\left(n\,-\,1\right)\,\rho\right)\mathrm{\mathbf{I}}_{n}\,-\,\rho\,\mathbf{J}_{n}}{\left(1\,-\,\rho\right)\left(1\,+\,\left(n\,-\,1\right)\,\rho\right)}.\label{eq:ConstantCorrelationShermanMorrison}
\end{align}


\subsection{The Inverse of a Covariance Matrix}

Sometime, a closed form inverse of a matrix yields useful insights
into the nature of its entries, and \citet{Stevens1998} provides
a particularly good example of this. Assume that we have a set of
$n$ securities. The covariance matrix of their returns can be partitioned
as

\begin{align}
\mathrm{\mathbf{C}}_{n} & =\left[\begin{array}{c|c}
\mathrm{\sigma}_{1}^{2}\, & \,\mathbf{\Sigma}_{12}\\
\hline \mathbf{\Sigma}_{21}\, & \,\mathrm{\mathbf{C}}_{n\,-\,1}
\end{array}\right],\label{eq:GuyStevensCovMat}
\end{align}

where $\sigma_{1}^{2}$ is the variance of security 1, $\mathbf{\Sigma}_{12}$
is a $1\times(n-1)$ vector that contains the covariances of security
1 with the remaining $n-1$ securities, $\mathbf{\Sigma}_{21}=\mathbf{\Sigma}_{12}^{\prime}$,
and $\mathrm{\mathbf{C}}_{n\,-\,1}$ is the covariance matrix of the
remaining $n-1$ securities. We can partition $\mathrm{\mathbf{C}}_{n}^{-1}$
in the same way we partition $\mathrm{\mathbf{C}}_{n}$:

\begin{equation}
\mathrm{\mathbf{C}}_{n}^{-1}=\left[\begin{array}{c|c}
a_{11}\, & \,\mathbf{\Lambda}{}_{12}\\
\hline \mathrm{\mathbf{\Lambda}}_{21}\, & \,\mathrm{\mathbf{\Lambda}}_{22}
\end{array}\right].\label{eq:GuyStevensInverse1}
\end{equation}

As 

\begin{align}
\mathrm{\mathbf{C}}_{n}^{-1}\mathrm{\mathbf{C}}_{n} & =\mathrm{\mathbf{I}}_{n},\nonumber \\
 & =\left[\begin{array}{c|c}
1\, & \,\mathbf{0}\\
\hline \mathbf{0}\, & \,\mathrm{\mathbf{I}}_{n\,-\,1}
\end{array}\right]\label{eq:GuyStevensInverse2}
\end{align}
we can write

\begin{subequations}
\begin{align}
a_{11}\,\mathrm{\sigma}_{1}^{2}\,+\,\mathbf{\Lambda}_{12}\,\mathbf{\Sigma}_{21} & =1\label{eq:GuyStevensInverse3}\\[0.5em]
a_{11}\,\mathbf{\Sigma}_{12}\,+\,\mathbf{\Lambda}_{12}\,\mathrm{\mathbf{C}}_{n-1} & =\mathbf{0}.\label{eq:GuyStevensInverse4}
\end{align}
\end{subequations}

On solving \eqref{eq:GuyStevensInverse4} for $\mathbf{\Lambda}_{12}$,
substituting into \eqref{eq:GuyStevensInverse3} and simplifying,
we get

\begin{align}
a_{11} & =\frac{1}{\mathrm{\sigma}_{1}^{2}\,-\,\mathbf{\Sigma}_{12}\,\mathrm{\mathbf{C}}_{n\,-\,1}^{-1}\,\mathbf{\Sigma}_{21}}\nonumber \\
 & =\frac{1}{\mathrm{\sigma}_{1}^{2}\,\left(1\,-\,R_{1}^{2}\right)}\nonumber \\
 & =\frac{1}{\mathrm{\sigma}_{\epsilon}^{2}},\label{eq:GuyStevensA11}
\end{align}

where $R_{1}^{2}$ is the fraction of the variance of security 1 that
is explained by securities $2,3,\cdots,n$ when its returns are regressed
against theirs, and $\mathrm{\sigma}_{\epsilon}^{2}$ is the variance
of the regression residuals. These are well known results in the theory
of least squares, and are discussed in \citet{Stevens1997} and \citet{Hayashi2000}.
Next, we write

\begin{align}
\mathbf{\Lambda}_{12} & =-a_{11}\,\mathbf{\Sigma}_{12}\,\mathrm{\mathbf{C}}_{n\,-\,1}^{-1}\nonumber \\
 & =\frac{-\mathbf{\Sigma}_{12}\,\mathrm{\mathbf{C}}_{n\,-\,1}^{-1}}{\mathrm{\sigma}_{1}^{2}\,\left(1\,-\,R_{1}^{2}\right)}\nonumber \\
 & =\frac{-\mathrm{\boldsymbol{\beta}_{1}}}{\mathrm{\sigma}_{1}^{2}\,\left(1\,-\,R_{1}^{2}\right)},\label{eq:GuyStevensBeta1}
\end{align}

where $\boldsymbol{\beta}_{1}$ is the vector of regression coefficients
obtained when security 1's returns are regressed against those of
securities $1,\cdots,\,i-1,\,i+1,\cdots,\,n$. 

But since our choice of security 1 is arbitrary, all the rows must
have exactly the same form, and the inverse of the covariance matrix
must be

\begin{equation}
\mathrm{\mathbf{C}}_{n}^{-1}=\left[\begin{array}{cccc}
\frac{1}{\mathrm{\sigma}_{1}^{2}\left(1\,-\,R_{1}^{2}\right)}\, & \,\frac{-\beta_{12}}{\mathrm{\sigma}_{1}^{2}\left(1\,-\,R_{1}^{2}\right)}\, & \,\cdots\, & \,\frac{-\beta_{1n}}{\mathrm{\sigma}_{1}^{2}\left(1\,-\,R_{1}^{2}\right)}\\
\frac{-\beta_{21}}{\mathrm{\sigma}_{2}^{2}\left(1\,-\,R_{2}^{2}\right)}\, & \,\frac{1}{\mathrm{\sigma}_{2}^{2}\left(1\,-\,R_{2}^{2}\right)}\, &  & \,\frac{-\beta_{2n}}{\mathrm{\sigma}_{2}^{2}\left(1\,-\,R_{2}^{2}\right)}\\
\vdots & \vdots & \,\ddots\, & \vdots\\
\frac{-\beta_{n1}}{\mathrm{\sigma}_{n}^{2}\left(1\,-\,R_{n}^{2}\right)}\, & \,\frac{-\beta_{n2}}{\mathrm{\sigma}_{n}^{2}\left(1\,-\,R_{n}^{2}\right)}\, &  & \,\frac{1}{\mathrm{\sigma}_{n}^{2}\left(1\,-\,R_{n}^{2}\right)}
\end{array}\right].\label{eq:GuyStevensInverse5}
\end{equation}

Even though we would not ever invert a covariance matrix in this way
(both the data requirements and the computational effort are huge,
and the procedure is not numerically stable), the insights it affords
are valuable, and helps clarify our thinking about the structure of
the inverse of a covariance matrix.

\subsection{The Mahanalobis Distance}

It sometimes proves useful to measure the distance between two random
vectors that are drawn from the same \emph{joint} distribution, but
whose \emph{marginal} distributions are very different. For example,
vectors $\mathrm{\mathbf{u}}$ and $\mathrm{\mathbf{v}}$ may contain
the returns of all the stocks in the S\&P 500 on two adjacent days,
but the marginal distribution of their first entries (a volatile small
cap biotech firm) may be very different than that of their last entry
(a steadily growing consumer goods company). In such cases, it proves
useful to first express the distance between each component of the
two vectors in standard deviation units, and to then compute a standardized
distance between them using a variant of \eqref{eq:DistanceBetweenVectors}.
In the special case when all the components of $\mathrm{\mathbf{u}}$
and $\mathrm{\mathbf{v}}$ are independent, this gives us

\begin{equation}
d_{\mathrm{\mathbf{u}},\mathrm{\mathbf{v}}}=\sqrt{\left(\frac{u_{1}\,-\,v_{1}}{\sigma_{1}}\right)^{2}\,+\,\cdots\,+\,\left(\frac{u_{n}\,-\,v_{n}}{\sigma_{n}}\right)^{2}}.\label{eq:DistanceScaledBySigma}
\end{equation}

More generally, the \emph{Mahalanobis distance} takes into account
the covariances between the components of the vectors, and is defined
as

\begin{equation}
d_{\mathrm{\mathbf{u}},\mathrm{\mathbf{v}}}^{M}=\sqrt{\left(\mathrm{\mathbf{u}}-\mathrm{\mathbf{v}}\right)^{\prime}\,\mathrm{\mathbf{C}}^{-1}\,\left(\mathrm{\mathbf{u}}-\mathrm{\mathbf{v}}\right)},\label{eq:MahanalobisDistance}
\end{equation}

where $C$ is the covariance matrix of the components of $\mathrm{\mathbf{u}}$
(and $\mathrm{\mathbf{v}}$). If the covariance matrix is diagonal
(i.e. if all the components of a vector are uncorrelated with each
other), \eqref{eq:MahanalobisDistance} reduces to \eqref{eq:DistanceScaledBySigma}.

\subsection{Eigenvalues and Eigenvectors}

In \eqref{eq:AxEQb}, a matrix $\mathrm{\mathbf{A}}$ maps a non--zero
vector $\mathrm{\mathbf{x}}$ to another vector $\mathrm{\mathbf{b}}$
which may lie at a great distance from $\mathrm{\mathbf{x}}$. But
not all matrices push vectors far away: the identity matrix $\mathrm{\mathbf{I}}_{n}$,
for example, maps every vector to itself, and a scalar multiple of
$\mathrm{\mathbf{I}}_{n}$ maps a vector to the same scalar multiple
of itself. This property is not unique to $\mathrm{\mathbf{I}}_{n}$
-- all square matrices map certain special vectors to scaled versions
of themselves, i.e

\begin{equation}
\mathrm{\mathbf{A}\,\mathrm{\mathbf{e}}=\lambda\,\mathrm{\mathbf{e}}.}\label{eq:Eigenvector1}
\end{equation}

The vector $\mathrm{\mathbf{e}}$ is known as an \emph{eigenvector}
of $\mathrm{\mathbf{A}}$, the scalar $\lambda$ is its corresponding
\emph{eigenvalue}, and they are collectively referred to as an \emph{eigenpair}.
Note that $\mathrm{\mathbf{e}}$ is not unique -- any scalar multiple
of $\mathrm{\mathbf{e}}$ is also an eigenvector of $\mathrm{\mathbf{A}}$.
Often, an eigenvector is scaled so that its norm is 1, and the resulting
eigenvector is known as a \emph{unit eigenvector} or a \emph{normalized
eigenvector}.

If we rewrite \eqref{eq:Eigenvector1} as

\begin{equation}
\left(\mathbf{A}\,-\,\lambda\,\mathrm{\mathbf{I}}\right)\mathrm{\mathbf{e}}=\mathrm{\mathbf{0},}\label{eq:Eigenvector2}
\end{equation}

the eigenvalues of $\mathbf{A}$ can be seen to be those values of
$\lambda$ for which $\mathbf{A}-\lambda\mathrm{\mathbf{I}}$ is singular,
or non--invertible, and therefore has a null space, or \emph{eigenspace};
any vector $\mathbf{e}$ that lies in the eigenspace associated with
a particular eigenvalue is an eigenvector of $\mathrm{\mathbf{A}}$.
It follows from the definition of invertibility that 
\begin{equation}
det\left(\mathbf{A}\,-\,\lambda\,\mathrm{\mathbf{I}}\right)=0,\label{eq:Eigenenvector3}
\end{equation}

Let $\lambda_{i},1\leqslant i\leqslant n$ be an eigenvalue of an
$n\times n$ matrix $\mathrm{\mathbf{A}}$. Then

\begin{align}
\sum_{i\,=\,1}^{n}\lambda_{i} & =\sum_{i\,=\,1}^{n}a_{ii}\nonumber \\
 & \triangleq trace\left(\mathrm{\mathbf{A}}\right),\label{eq:SumOfEigenvalues}
\end{align}

and

\begin{equation}
\prod_{i\,=\,1}^{n}\lambda_{i}=det\left(\mathrm{\mathbf{A}}\right).\label{eq:ProductOfEigenvalues}
\end{equation}

The set of eigenvalues of a matrix $\mathrm{\mathbf{A}}$ is known
as its \emph{spectrum}, and its spectral radius, $\rho\left(\mathbf{A}\right)$,
is defined to be 

\begin{equation}
\rho\left(\mathbf{A}\right)=\max_{i}\left|\lambda_{i}\right|,\label{eq:SpectralRadius1}
\end{equation}

and it can be shown that

\begin{equation}
\rho\left(\mathbf{A}\right)\leqslant\min\left(\max_{i}\sum_{j\,=\,1}^{n}\left|a_{ij}\right|,\ \max_{j}\sum_{i\,=\,1}^{n}\left|a_{ij}\right|\right).\label{eq:SpectralRadius2}
\end{equation}

A few useful facts about the eigenpairs $\lambda_{i}$ and $\mathrm{\mathbf{e}}_{i}$
of an $n\times n$ matrix $\mathbf{A}$ follow:
\begin{enumerate}
\item If $\mathbf{A}$ is invertible, $\nicefrac{1}{\lambda_{i}}$ and $\mathrm{\mathbf{e}}_{i}$
are an eigenpair of $\mathbf{A}^{-1}.$ 
\item The eigenvalues (but not necessarily the eigenvectors) of $\mathrm{\mathbf{A}^{\prime}}$
are identical to the eigenvalues of $\mathrm{\mathbf{A}}.$
\item $\lambda_{i}^{k}$ and $\mathrm{\mathbf{e}}_{i}$ are an eigenpair
of $\mathbf{A}^{k},\ k>1$.
\item If $\mathbf{A}$ is triangular or diagonal, then $\lambda_{i}=a_{ii},\ 1\leqslant i\leqslant n$
.
\item $\lambda_{i}\,-\,k$ and $\mathrm{\mathbf{e}}_{i}$ are an eigenpair
of $\mathbf{A\,-\,}k\,\mathrm{\mathbf{I}}_{n}$.
\item If $\mathbf{B}$ is an invertible $n\times n$ matrix, the eigenvalues
(but not necessarily the eigenvectors) of $\mathbf{B}^{-1}\,\mathbf{A}\,\mathbf{\mathbf{B}}$
are identical to those of $\mathbf{A}$.
\item If the eigenvalues of $\mathbf{A}$ are distinct, or if $\mathbf{A}$
is symmetric and has only real entries (covariance and correlation
matrices are good examples of this second, weaker, condition), the
matrix $\mathbf{P}$ formed by stacking its normalized (i.e. of length
1) eigenvectors side by side diagonalizes $\mathbf{A}$, i.e.
\end{enumerate}
\begin{equation}
\mathbf{P}^{-1}\,\mathbf{A}\,\mathbf{\mathbf{P}}=\mathbf{D}_{\lambda}\label{eq:EigenvaluesDiagonalizeA}
\end{equation}

where

\begin{equation}
\mathbf{D}_{\lambda}=\left[\begin{array}{cccc}
\lambda_{1} & 0 & \cdots & 0\\
0 & \ddots &  & \vdots\\
\vdots &  & \ddots & 0\\
0 & \cdots & 0 & \lambda_{n}
\end{array}\right].\label{eq:Dlambda}
\end{equation}

Equation \eqref{eq:EigenvaluesDiagonalizeA} is often written as

\begin{equation}
\mathbf{\mathbf{A}=\mathbf{P}\,\mathbf{D}_{\lambda}\,\mathbf{P}^{-1}},\label{eq:SpectralRepresentation}
\end{equation}

where it is called the \emph{eigendecomposition}, or the \emph{spectral}
\emph{representation}, of $\mathrm{\mathbf{A}}$. 

It can further be shown that when the eigenvalues of $\mathbf{A}$
are distinct, or if $\mathbf{A}$ is symmetric and has only real entries, 

\[
\mathbf{P}^{-1}=\mathbf{P}^{\prime},
\]

so that \eqref{eq:SpectralRepresentation} can usefully be rewritten
as

\begin{align}
\mathbf{\mathbf{A}} & =\mathbf{P}\,\mathbf{D}_{\lambda}\,\mathbf{P}^{\prime}\nonumber \\
 & =\sum_{i=1}^{n}\lambda_{i}\,\mathrm{\mathbf{p}}_{i}\,\mathrm{\mathbf{p}}_{i}^{\prime},\label{eq:SpectralRepresentation1}
\end{align}

where $\mathrm{\mathbf{p}}_{i}$ is the $i^{th}$ column of $\mathbf{P}$. 

Clearly, $\mathrm{\mathbf{A}}$ is invertible only if $\mathbf{D}_{\lambda}$
is invertible, and $\mathbf{D}_{\lambda}$ in turn is invertible if
and only if all its eigenvalues are non--zero. Direct multiplication
shows that

\begin{align}
\mathrm{\mathbf{A}}^{k} & =\mathbf{P}\,\mathbf{D}_{\lambda}^{k}\,\mathbf{P}^{-1}\nonumber \\
 & =\mathbf{P}\,\mathbf{D}_{\lambda}^{k}\,\mathbf{P}^{\prime}\nonumber \\
 & =\sum_{i\,=\,1}^{n}\lambda_{i}^{k}\,\mathrm{\mathbf{p}}_{i}\,\mathrm{\mathbf{p}}_{i}^{\prime},\ k\geqslant2,\label{eq:SpectralRepresentationOfPowers}
\end{align}

under the above conditions, and it is also true that

\begin{align}
\exp\left(\mathrm{\mathbf{A}}\right) & =\mathbf{P}\,\exp\left(\mathbf{D}_{\lambda}\right)\,\mathbf{P}^{-1}\nonumber \\
 & =\mathbf{P}\,\exp\left(\mathbf{D}_{\lambda}\right)\,\mathbf{P}^{\prime},\label{eq:SpectralRepresentationOfExpA}
\end{align}

where 

\begin{equation}
\exp\left(\mathrm{\mathbf{A}}\right)\triangleq\mathbf{I}_{n}\,+\,\sum_{i\,>\,0}\frac{1}{i!}\mathrm{\mathbf{A}}^{i}\label{eq:MatrixExponential}
\end{equation}
 

is known as the \emph{matrix exponential}. 

It follows that

\begin{equation}
\exp\left(\mathrm{\mathbf{0}}_{n}\right)=\mathbf{I}_{n}.\label{eq:MatrixExponential0}
\end{equation}

In the special case when $\mathbf{A}$ and $\mathbf{B}$ commute,
i.e. when $\mathbf{A}\,\mathbf{B}=\mathbf{B}\,\mathbf{A}$,

\begin{align}
\exp\left(\mathrm{\mathbf{A}}\right)\exp\left(\mathrm{\mathbf{B}}\right) & =\exp\left(\mathrm{\mathbf{B}}\right)\,\exp\left(\mathrm{\mathbf{A}}\right)\nonumber \\
 & =\exp\left(\mathrm{\mathbf{A}}\,+\,\mathrm{\mathbf{B}}\right).\label{eq:MatrixExponentialProduct}
\end{align}

In particular, as $\mathrm{\mathbf{A}}$ commutes with itself, and
as $\mathrm{\mathbf{A}\,+}\,\left(-\mathrm{\mathbf{A}}\right)=\mathbf{0}_{n}$

\begin{equation}
\exp\left(\mathrm{\mathbf{A}}\right)^{-1}=\exp\left(-\mathrm{\mathbf{A}}\right).\label{eq:MatrixExponentialInverse}
\end{equation}

Equation \eqref{eq:EigenvaluesDiagonalizeA} can be generalized: two
matrices, $\mathrm{\mathbf{A}}$ and $\mathrm{\mathbf{B}}$, are said
to be \emph{similar} if there exists an invertible matrix $\mathrm{\mathbf{E}}$
such that 

\begin{equation}
\mathbf{E}^{-1}\,\mathbf{A}\,\mathbf{\mathbf{E}}=\mathbf{B}.\label{eq:SimilarMatrices1}
\end{equation}

Similar matrices have identical eigenvalues, as

\begin{align}
det\left(\mathbf{B}\,-\,\lambda\mathrm{\mathbf{I}}\right) & =det\left(\mathbf{\mathbf{E}^{-1}}\left(\mathbf{A}\,-\,\lambda\,\mathrm{\mathbf{I}}\right)\mathbf{\mathbf{E}}\right)\nonumber \\
 & =det\left(\mathbf{\mathbf{E}^{-1}}\right)det\left(\mathbf{A}\,-\,\lambda\,\mathrm{\mathbf{I}}\right)det\left(\mathbf{\mathbf{E}}\right)\nonumber \\
 & =det\left(\mathbf{A}\,-\,\lambda\,\mathrm{\mathbf{I}}\right).\label{eq:SimilarMatrices2}
\end{align}


\subsection{Positive Definite Matrices}

A matrix $\mathbf{A}$ is said to be \emph{positive definite} if 

\begin{equation}
\mathrm{\mathbf{x}}^{\prime}\,\mathrm{\mathbf{A\,x}}>\mathrm{\mathbf{0}}\ \forall\mathrm{\mathbf{x}}\ne\mathrm{\mathbf{0}}.\label{eq:PositiveDefinite}
\end{equation}
An equivalent condition for positive definiteness is that all the
eigenvalues of $\mathbf{A}$ are strictly positive. If we relax the
inequality slightly, i.e. if

\begin{equation}
\mathrm{\mathbf{x}}^{\prime}\,\mathrm{\mathbf{A\,x}}\geqslant\mathrm{\mathbf{0}}\ \forall\mathrm{\mathbf{x}}\ne\mathrm{\mathbf{0}},\label{eq:PositiveSemiDefinite}
\end{equation}

$\mathbf{A}$ is said to be \emph{positive semi--definite}. 

Negating each entry in a positive definitive (or positive semi--definite)
matrix results in a \emph{negative definite} (or \emph{negative semi--definite})
matrix -- the inequalities in \eqref{eq:PositiveDefinite} and \eqref{eq:PositiveSemiDefinite}
are reversed.

The Gramian matrix $\mathbf{A}^{\prime}\mathrm{\mathbf{A}}$ is non--negative
definite, as

\begin{align}
\mathrm{\mathbf{x}}^{\prime}\,\mathrm{\mathbf{\mathbf{A}^{\prime}\,\mathrm{\mathbf{A}\,}x}} & =\left(\mathbf{\mathbf{A\,}x}\right)^{\prime}\left(\mathbf{\mathbf{A\,}x}\right)\nonumber \\
 & \geqslant\mathrm{\mathbf{0}\ \forall\mathrm{\mathbf{x}}}.\label{eq:GramianNonNegativeDefinite}
\end{align}

Furthermore, the rank of the Gramian is identical to the rank of $\mathbf{\mathbf{A}}$.

If $\mathbf{A}$ is symmetric and positive semi--definite, its eigenvalues
are non--negative. If it is symmetric and positive definite, its
eigenvalues are strictly positive.

Positive definite matrices are of great importance in portfolio management
because they are used to compute the variance of a portfolio from
the variances and covariances of its constituent securities (the variance
of a portfolio is, by definition, non--negative), and also find use
when identifying extrema of multivariate functions. 

As a positive definite covariance matrix is symmetric and has only
real entries, \eqref{eq:EigenvaluesDiagonalizeA}--\eqref{eq:SpectralRepresentation1}
apply, and we can write

\begin{equation}
\mathbf{P}^{\prime}\,\mathbf{A}\,\mathbf{P}=\mathbf{D}_{\lambda},\label{eq:PDMatrixDiagonalization}
\end{equation}

where $\mathbf{P}$ is an orthonormal matrix whose columns are the
normalized eigenvalues of $\mathbf{A}$ and $\mathbf{D}_{\lambda}$
is a diagonal matrix of eigenvalues -- see \eqref{eq:Dlambda}. If
we define

\begin{equation}
\sqrt{\mathbf{D}_{\lambda}^{-1}}=\left[\begin{array}{cccc}
\nicefrac{1}{\sqrt{\lambda_{1}}} & 0 & \cdots & 0\\
0 & \ddots &  & \vdots\\
\vdots &  & \ddots & 0\\
0 & \cdots & 0 & \nicefrac{1}{\sqrt{\lambda_{n}}}
\end{array}\right]\label{eq:SqrtDlambdaInv}
\end{equation}

and

\begin{equation}
\mathbf{Q}=\mathbf{P}\,\sqrt{\mathbf{D}_{\lambda}^{-1}},\label{eq:QFromP}
\end{equation}

we can write

\begin{equation}
\mathbf{Q}^{\prime}\,\mathbf{A}\,\mathbf{Q}=\mathbf{I}_{n}.\label{eq:CovarianceMatrixInverse}
\end{equation}

Using straightforward properties of Cholesky factorization and products
of orthogonal matrices, it can be shown that $\mathbf{Q}$ diagonalizes
\emph{any} positive definite covariance matrix, say $\mathbf{B}$,
i.e.

\begin{equation}
\mathbf{Q}^{\prime}\,\mathbf{B}\,\mathbf{Q}=\mathbf{D}_{1},\label{eq:PdiagonalizesB}
\end{equation}

where $\mathbf{D}_{1}$ is a diagonal matrix, though its entries are
not the eigenvalues of $\mathbf{B}$. The matrices $\mathbf{A}$ and
$\mathbf{B}$ are said to be simultaneously diagonalizable.


\subsection{QR Decomposition and the Gram--Schmidt Procedure\label{subsec:QR_Decomposition}}

Assume that we have a set of vectors $\mathrm{\mathbf{q}}_{1},\,\mathrm{\mathbf{q}}_{2,}\,\cdots\,,\mathrm{\mathbf{q}}_{n}$
which form an orthonormal basis\footnote{An orthogonal basis can be transformed into an orthonormal basis by
dividing each vector by its norm.} for an $n$--dimensional vector space $\mathcal{{S}}$ . From the
definition of orthonormality,

\begin{equation}
\left\langle \mathrm{\mathbf{q}}_{i},\,\mathrm{\mathbf{q}}_{j}\right\rangle =\begin{cases}
1 & i=j\\
0 & i\ne j.
\end{cases}\label{eq:Orthonormality}
\end{equation}

If the $\mathrm{\mathbf{q}}_{i}$ are vectors in $n$--dimensional
Euclidean space, we can stack them side by side to form an $n\times n$
matrix

\begin{equation}
\mathrm{\mathbf{Q}}=\left[\mathrm{\mathbf{q}}_{1}\,\cdots\,\mathrm{\mathbf{q}}_{n}\right],\label{eq:Qmatrix}
\end{equation}

and it is easily seen that 

\begin{align}
\mathrm{\mathbf{Q}}^{\prime}\,\mathrm{\mathbf{Q}} & =\mathrm{\mathbf{I}}_{n}.\label{eq:QtransposeQequalsIn}
\end{align}

We can post--multiply both sides of \eqref{eq:QtransposeQequalsIn}
by $\mathrm{\mathbf{Q}}^{-1}$ to obtain

\begin{equation}
\mathrm{\mathbf{Q}}^{-1}=\mathrm{\mathbf{Q}}^{\prime},\label{eq:QinverseEqualsQtranspose}
\end{equation}

but it is important to note that \eqref{eq:QtransposeQequalsIn} (but
not \eqref{eq:QinverseEqualsQtranspose}) holds even if $\mathrm{\mathbf{Q}}$
is rectangular, i.e. if the $\mathrm{\mathbf{q}}_{i}$ are vectors
in an $m$--dimensional Euclidean space, for $m>n.$ Also, as $\mathrm{\mathbf{I}}_{n}^{\prime}=\mathrm{\mathbf{I}}_{n}$,
it follows that

\begin{equation}
\mathrm{\mathbf{Q}}^{\prime}\,\mathrm{\mathbf{Q}}=\mathrm{\mathbf{Q}}\,\mathrm{\mathbf{Q}}^{\prime}.\label{eq:QtransposeQ_2}
\end{equation}

Remarkably, \emph{any} matrix whose columns are linearly independent
can be decomposed into a product of a orthonormal matrix and an invertible
upper triangular matrix\footnote{Normally, an upper triangular matrix would be represented by $\mathrm{\mathbf{U}}$,
but it has become the norm to represent it in this particular factorization
by $\mathrm{\mathbf{R}}$, possibly because all its entries lie to
the right of the principal diagonal.}, i.e.

\begin{equation}
\mathrm{\mathbf{A}}=\mathrm{\mathbf{Q}}\,\mathrm{\mathbf{R}}\label{eq:QRdecomposition}
\end{equation}

where $\mathrm{\mathbf{Q}}$ is orthonormal and $\mathrm{\mathbf{R}}$
is invertible and upper triangular.

The \emph{Gram--Schmidt} procedure is a constructive algorithm for
creating $\mathrm{\mathbf{Q}}$ from any full rank matrix $\mathrm{\mathbf{A}}$. 

Let

\begin{equation}
\mathrm{\mathbf{A}}=\left[\mathrm{\mathbf{a}}_{1}\,\cdots\,\mathrm{\mathbf{a}}_{n}\right].\label{eq:Amatrix}
\end{equation}

be an $n\times n$ matrix of full rank.

If we divide each column of $\mathrm{\mathbf{A}}$ by its norm, we
get an equivalent\footnote{Equivalent in the sense that the column space of $\mathrm{\mathbf{A}}_{1}$
is identical to the column space of $\mathrm{\mathbf{A}}$.} matrix $\mathrm{\mathbf{A}}_{1}$, whose columns all have norm 1. 

\begin{align}
\mathrm{\mathbf{A}}_{1} & =\left[\frac{\mathrm{\mathbf{a}}_{1}}{\left\Vert \mathrm{\mathbf{a}}_{1}\right\Vert }\,\cdots\,\frac{\mathrm{\mathbf{a}}_{n}}{\left\Vert \mathrm{\mathbf{a}}_{n}\right\Vert }\right]\nonumber \\
 & \triangleq\left[\mathrm{\mathbf{a}}_{1}^{1}\,\cdots\,\mathrm{\mathbf{a}}_{n}^{1}\right].\label{eq:A1matrix}
\end{align}

To construct $\mathbf{Q}$ using the Gram--Schmidt procedure, we
first set $\mathrm{\mathbf{q}}_{1}=\mathrm{\mathbf{a}}_{1}^{1}$,
and then iteratively subtract the projections of $\mathrm{\mathbf{a}}_{i}^{1}$
for $i=2,\,3,\,\cdots\,,\,n$ onto $\mathrm{\mathbf{q}}_{1},\,\cdots\,,\,\mathrm{\mathbf{q}}_{i-1}$
from itself, and finally normalize the resulting vector to create
a new vector that is orthonormal to all the preceding vectors. Formally,

\begin{equation}
\begin{array}{cl}
\mathrm{\mathbf{q}}_{1} & =\mathrm{\mathbf{a}}_{1}^{1}\\
\mathrm{\mathbf{q}}_{2} & =\frac{\mathbf{\mathrm{\mathbf{a}}}_{2}^{1}\,-\,\left\langle \mathbf{\mathrm{\mathbf{a}}}_{2}^{1},\,\mathrm{\mathbf{q}}_{1}\right\rangle \mathrm{\mathbf{q}}_{1}}{\left\Vert \mathbf{\mathrm{\mathbf{a}}}_{2}^{1}\,-\,\left\langle \mathbf{\mathrm{\mathbf{a}}}_{2}^{1},\,\mathrm{\mathbf{q}}_{1}\right\rangle \mathrm{\mathbf{q}}_{1}\right\Vert }\\
 & \vdots\\
\mathrm{\mathbf{q}}_{n} & =\frac{\mathbf{\mathrm{\mathbf{a}}}_{n}^{1}\,-\,\sum_{i\,=\,1}^{n\,-\,1}\left\langle \mathbf{\mathrm{\mathbf{a}}}_{n}^{1},\,\mathrm{\mathbf{q}}_{i}\right\rangle \mathrm{\mathbf{q}}_{i}}{\left\Vert \mathbf{\mathrm{\mathbf{a}}}_{n}^{1}\,-\,\sum_{i\,=\,1}^{n\,-\,1}\left\langle \mathbf{\mathrm{\mathbf{a}}}_{n}^{1},\,\mathrm{\mathbf{q}}_{i}\right\rangle \mathrm{\mathbf{q}}_{i}\right\Vert }
\end{array}\label{eq:GramSchmidt}
\end{equation}

It is easy to check that the $\mathrm{\mathbf{q}}_{i}$ satisfy \eqref{eq:Orthonormality},
and that the column space of $\mathrm{\mathbf{Q}}$ is identical to
the column space of $\mathrm{\mathbf{A}}$. Keep in mind that the
Gram--Schmidt procedure does not return a unique orthonormal matrix
-- if, for example, the columns of $\mathrm{\mathbf{A}}$ are resequenced,
it will return a different orthonormal matrix.

Even though the Gram--Schmidt procedure makes clear the conceptual
basis of orthonormalization, it is rarely used in practice on account
of its numerically instability when the columns of $\mathrm{\mathbf{A}}$
are almost collinear. And though its numerical stability can be improved
by resequencing its computations (see \citet{Gentle2017} for details),
\emph{Householder transformations} and \emph{Givens rotations} are
more widely used in practice to create $\mathrm{\mathbf{Q}}\mathbf{R}$
decompositions, which see widespread use in the solution of \eqref{eq:AxEQb}. 

The fact that the transpose of an orthonormal matrix is its inverse
leads to a subtly different notion of similarity: two matrices, $\mathrm{\mathbf{A}}$
and $\mathrm{\mathbf{B}}$, are said to be \emph{orthogonally similar}
if there exists an orthonormal matrix $\mathrm{\mathbf{F}}$ such
that 

\begin{equation}
\mathbf{F}^{\prime}\,\mathbf{A\,}\mathbf{\mathbf{F}}=\mathbf{B}.\label{eq:SimilarMatrices3}
\end{equation}

If $\mathbf{B}$ is upper triangular, then $\mathbf{F}\,\mathbf{B\,\mathbf{F}^{\prime}}$is
known as the \emph{Schur factorization} of $\mathbf{A}$.

\subsection{The Singular Value Decomposition}

In the special case when $\mathbf{A}$ is square and symmetric, it
can be represented as the Gramian of some $n\times m$ matrix $\mathbf{B}$
with rank $r\leqslant\min\left(m,n\right)$ (think of the columns
of $\mathbf{B}$ as time series of returns of securities, some of
which may admit a riskless arbitrage), i.e.,

\begin{equation}
\mathbf{\mathbf{A}}=\mathbf{B}^{\prime}\,\mathbf{B}.\label{eq:GramianOfPD}
\end{equation}

The rank of $\mathbf{\mathbf{A}}$ is also $r$, and it possesses
a set of orthonormal eigenvectors that diagonalize it, i.e. its spectral
representation is

\begin{align}
\mathbf{\mathbf{A}} & =\mathbf{Q}\,\mathbf{D}_{\lambda}\,\mathbf{Q}^{\prime},\label{eq:PDSpectralRepresentation}
\end{align}

where $\mathbf{Q}$ is an orthonormal matrix whose columns are the
normalized eigenvectors of $\mathbf{A}$, and $\mathbf{D}_{\lambda}$
is a diagonal matrix whose entries are the eigenvalues of $\mathbf{A}$.

The \emph{Singular Value Decomposition} (SVD) generalizes and extends
\eqref{eq:PDSpectralRepresentation} to \emph{all} matrices, and diagonalizes
$\mathbf{\mathbf{A}}_{n\times m}$ using its \emph{singular vectors}
instead of its eigenvalues. The resulting decomposition represents
any $n\times m$ matrix as a product of three matrices

\begin{equation}
\mathbf{\mathbf{A}}_{n\,\times\,m}=\mathbf{U}_{n\,\times\,n}\,\mathbf{\Sigma}_{n\,\times\,m}\,\mathbf{V}_{m\,\times\,m}^{\prime},\label{eq:SVD}
\end{equation}

where $\mathbf{U}$ and $\mathbf{V}$ are orthonormal matrices, and
$\mathbf{\mathbf{\Sigma}}$ is diagonal in the following sense:

\begin{equation}
\sigma_{ij}=\begin{cases}
\sigma_{i}>0 & 1\leqslant i=j\leqslant\mathrm{rank}\left(\mathbf{A}\right)\\
0 & \mathrm{otherwise}.
\end{cases}\label{eq:DiagMatrixInSVD}
\end{equation}

The $\sigma_{i}$ are known as \emph{singular values}, and are the
square roots of the non-zero eigenvalues of $\mathbf{A}^{\prime}\,\mathbf{A}$.
Correspondingly, the columns of $\mathbf{U}$ are the orthonormal
eigenvectors of $\mathbf{A}\,\mathbf{A}^{\prime}$ and those of $\mathbf{V}$
are the orthonormal eigenvectors of $\mathbf{A}^{\prime}\,\mathbf{A}$.
They are known as the \emph{left} and \emph{right singular vectors},
respectively, of $\mathrm{\mathbf{A}}$. For symmetric, square, matrices,
\eqref{eq:SVD} reduces to \eqref{eq:PDSpectralRepresentation}.

Recognizing that $\mathbf{V}^{-1}=\mathbf{V}^{\prime},$we can rewrite
\eqref{eq:SVD} as

\begin{equation}
\mathbf{\mathbf{A}}_{n\,\times\,m}\,\mathbf{V}_{m\,\times\,m}=\mathbf{U}_{n\,\times\,n}\,\mathbf{\Sigma}_{n\,\times\,m},\label{eq:SVD1}
\end{equation}

so that $\mathbf{\mathbf{A}}$ maps the columns of $\mathbf{V}$ to
scaled versions of the columns of $\mathbf{U}$. Addtionally, we can
partition $\mathbf{U}$ and $\mathbf{V}$ into two adjacent block
matrices 

\begin{subequations}
\begin{align}
\mathbf{U} & =\left[\mathrm{\mathbf{U}}_{1}\ |\ \mathrm{\mathbf{U}}_{2}\right]=\left[\begin{array}{ccc}
\mathrm{\mathbf{u}}_{1} & \,\cdots & \ \mathrm{\mathbf{u}}_{r}\ |\ \mathrm{\mathbf{u}}_{r+1}\,\cdots\ \mathrm{\mathbf{u}}_{n}\end{array}\right]\label{eq:SVDu}\\[0.5em]
\mathbf{V} & =\left[\mathrm{\mathbf{V}}_{1}\ |\ \mathrm{\mathbf{V}}_{2}\right]=\left[\begin{array}{ccc}
\mathrm{\mathbf{v}}_{1} & \,\cdots & \ \mathrm{\mathbf{v}}_{r}\ |\ \mathrm{\mathbf{v}}_{r+1}\,\cdots\ \mathrm{\mathbf{v}}_{m}\end{array}\right],\label{eq:SVDv}
\end{align}
\end{subequations}

and it can be shown that

\[
\begin{array}{cl}
\mathbf{U}_{1} & =\left[\begin{array}{ccc}
\mathrm{\mathbf{u}}_{1} & \,\cdots\  & \mathrm{\mathbf{u}}_{r}\end{array}\right]\hspace{0.315cm}\mathrm{is\ an\ orthonormal\boldsymbol{\ }basis\ for\ the\ \mathit{\mathit{column}\ space}\ of\ \mathbf{\mathbf{A}}}\\
\mathbf{U}_{2} & =\left[\mathrm{\mathbf{u}}_{r\,+\,1}\,\cdots\ \mathrm{\mathbf{u}}_{n}\right]\hspace{0.12cm}\mathrm{is\ an\ orthonormal\boldsymbol{\ }basis\ for\ \mathbf{\mathcal{{N}}\left(\mathrm{\mathbf{A}}^{\prime}\right)},\ the\ \mathit{left\ null\ space}\ of}\ \mathrm{\mathbf{A}}\\
\mathbf{V}_{1} & =\left[\begin{array}{ccc}
\mathrm{\mathbf{v}}_{1} & \,\cdots\  & \mathrm{\mathbf{v}}_{r}\end{array}\right]\hspace{0.335cm}\mathrm{is\ an\ orthonormal\boldsymbol{\ }basis\ for\ the}\ row\ \mathit{space}\ \mathrm{of}\ \mathbf{\mathbf{A}}\\
\mathbf{V}_{2} & =\left[\mathrm{\mathbf{v}}_{r\,+\,1}\,\cdots\ \mathrm{\mathbf{v}}_{m}\right]\ \mathrm{is\ an\ orthonormal\boldsymbol{\ }basis\ for\ \mathbf{\mathcal{{N}}\left(\mathrm{\mathbf{A}}\right)},\ the}\ \mathit{null\ space}\ \mathrm{of}\ \mathbf{A}.
\end{array}
\]

If we write out \eqref{eq:SVD} in terms of the left and right singular
vectors $\mathrm{\mathbf{u}}_{i}$ and $\mathrm{\mathbf{v}}_{i}$
contained in $\mathbf{U}_{1}$ and $\mathbf{V}_{1}$ (the orthonormal
bases for its column space and row space respectively), we see that
the SVD allows us to represent$\mathrm{\mathbf{A}}$ as a sum of rank
1 matrices:

\begin{equation}
\mathbf{\mathbf{A}}=\sum_{i\,=\,1}^{r}\sigma_{i}\,\mathrm{\mathbf{u}}_{i}\,\mathbf{v}_{i}^{\prime}.\label{eq:SVD2}
\end{equation}

If we further reorder the terms on the right hand side of \eqref{eq:SVD2}
in descending order of $\sigma_{i}$, we obtain a decomposition of
$\mathrm{\mathbf{A}}$ in which the terms are ranked in descending
order of importance, and we can obtain a sequence of increasingly
accurate approximations to $\mathrm{\mathbf{A}}$ by computing the
partial sums of the first $p<r$ terms of this reordered sequence. 

This is, in fact, the way in which the SVD is most often used --
a small, but important, subset of the terms in \eqref{eq:SVD2} is
used to construct an adequate approximation to $\mathrm{\mathbf{A}}$
that greatly simplifies computations using it. In many real life applications,
matrices with many millions of entries are well approximated by a
partial sum with ten or twenty terms, speeding up computations by
many orders of magnitude.

\section{Some Results on Functions}

\subsection{Derivatives of Functions}

Let $f\left(\mathrm{\mathbf{x}}\right)$ be a differentiable function
of $n$ variables $x_{1},\,x_{2},\cdots\,,\,x_{n}$.

The \emph{gradient} of $f\left(\mathrm{\mathbf{x}}\right),\ \nabla f\left(\mathrm{\mathbf{x}}\right)$,
is a vector of its partial derivatives, and is given by

\begin{equation}
\nabla f\left(\mathrm{\mathbf{x}}\right)=\left[\begin{array}{c}
\frac{\partial f\left(\mathrm{\mathbf{x}}\right)}{\partial x_{1}}\\
\vdots\\
\frac{\partial f\left(\mathrm{\mathbf{x}}\right)}{\partial x_{n}}
\end{array}\right].\label{eq:GradientFx}
\end{equation}

The gradient vector points in the direction in which $f$'s rate of
change is maximal. In the special case when 

\begin{align}
f\left(\mathrm{\mathbf{x}}\right) & =\sum_{i\,=\,1}^{n}a_{i}\,x_{i}\nonumber \\
 & =\mathrm{\mathbf{a}}\cdotp\mathrm{\mathbf{x}},\label{eq:fx_dot_product}
\end{align}

it is easily seen that

\begin{equation}
\nabla f\left(\mathrm{\mathbf{x}}\right)=\mathrm{\mathbf{a}}.\label{eq:GradientDotProduct}
\end{equation}

Next, if $f\left(\mathrm{\mathbf{x}}\right)$ is a \emph{quadratic
form}, i.e. if
\begin{align}
f\left(\mathrm{\mathbf{x}}\right) & =\sum_{i\,=\,1}^{n}\sum_{j\,=\,1}^{n}a_{ij}\,x_{i}\,x_{j}\nonumber \\
 & =\mathrm{\mathbf{x}}^{\prime}\,\mathbf{A}\,\mathrm{\mathbf{x}},\label{eq:fx_quadratic_form}
\end{align}

where the $a_{ij}$ are constants, it can be seen by differentiating
and gathering terms that

\begin{equation}
\nabla f\left(\mathrm{\mathbf{x}}\right)=2\,\mathbf{A}\,\mathrm{\mathbf{x}}.\label{eq:GradientQuadraticForm}
\end{equation}

If $f\left(\mathrm{\mathbf{x}}\right)$ is twice differentiable, its
Hessian, $\mathrm{\mathbf{H}}\left(f\left(\mathrm{\mathbf{x}}\right)\right)$,
is a symmetric square matrix of its second partial derivatives:

\begin{align}
\mathrm{\mathbf{H}}\left(f\left(\mathrm{\mathbf{x}}\right)\right) & =\frac{\partial^{2}f\left(\mathrm{\mathbf{x}}\right)}{\partial\mathrm{\mathbf{x}}\partial\mathrm{\mathbf{x}}^{\prime}}\nonumber \\
 & =\left[\begin{array}{cccc}
\frac{\partial^{2}f\left(\mathrm{\mathbf{x}}\right)}{\partial x_{1}^{2}}\, & \,\frac{\partial^{2}f\left(\mathrm{\mathbf{x}}\right)}{\partial x_{1}\,\partial x_{2}}\, & \cdots & \,\frac{\partial^{2}f\left(\mathrm{\mathbf{x}}\right)}{\partial x_{1}\,\partial x_{n}}\\
\frac{\partial^{2}f\left(\mathrm{\mathbf{x}}\right)}{\partial x_{2}\,\partial x_{1}}\, & \,\frac{\partial^{2}f\left(\mathrm{\mathbf{x}}\right)}{\partial x_{2}^{2}}\, & \cdots & \,\frac{\partial^{2}f\left(\mathrm{\mathbf{x}}\right)}{\partial x_{2}\,\partial x_{n}}\\
\vdots & \vdots & \ddots & \vdots\\
\frac{\partial^{2}f\left(\mathrm{\mathbf{x}}\right)}{\partial x_{n}\,\partial x_{1}}\, & \,\frac{\partial^{2}f\left(\mathrm{\mathbf{x}}\right)}{\partial x_{n}\,\partial x_{2}}\, & \cdots & \,\frac{\partial^{2}f\left(\mathrm{\mathbf{x}}\right)}{\partial x_{n}^{2}}
\end{array}\right]\label{eq:Hessian}
\end{align}

The Hessian matrix reflects the magnitude, as well as the principal
axes, of the local curvature of $f\left(\mathrm{\mathbf{x}}\right)$.

If 

\begin{equation}
\mathrm{\mathbf{g}}\left(x\right)=\left[\begin{array}{c}
g_{1}\left(x\right)\\
\vdots\\
g_{n}\left(x\right)
\end{array}\right]\label{eq:VectorOfFunctions}
\end{equation}

is a vector of $k$--times differentiable functions, $g_{1}\left(x\right),\,g_{2}\left(x\right),\,\cdots\,,\,g_{n}\left(x\right)$,
of a single variable $x$, then 

\begin{equation}
\frac{d^{k}\mathrm{\mathbf{g}}\left(x\right)}{dx^{k}}=\left[\begin{array}{c}
\frac{d^{k}g_{1}\left(x\right)}{dx^{k}}\\
\vdots\\
\frac{d^{k}g_{n}\left(x\right)}{dx^{k}}
\end{array}\right].\label{eq:VectorOfDerivatives}
\end{equation}

Equation \eqref{eq:VectorOfDerivatives} generalizes naturally to
matrices. If 

\begin{equation}
\mathrm{\mathbf{G}}\left(x\right)=\left[\begin{array}{ccc}
g_{11}\left(x\right) & \,\cdots\, & g_{1m}\left(x\right)\\
\vdots & \vdots & \vdots\\
g_{n1}\left(x\right) & \,\cdots\, & g_{nm}\left(x\right)
\end{array}\right]\label{eq:MatrixOfFunctions}
\end{equation}

is a matrix of $k$--times differentiable functions, $g_{11}\left(x\right),\,\cdots\,,\,g_{nm}\left(x\right)$,
of a single variable $x$, then

\begin{equation}
\frac{d^{k}\mathrm{\mathbf{G}}\left(x\right)}{dx^{k}}=\left[\begin{array}{ccc}
\frac{d^{k}g_{11}\left(x\right)}{dx^{k}} & \,\cdots\, & \frac{dg_{1m}^{k}\left(x\right)}{dx^{k}}\\
\vdots & \vdots & \vdots\\
\frac{d^{k}g_{n1}\left(x\right)}{dx^{k}} & \,\cdots\, & \frac{d^{k}g_{nm}\left(x\right)}{dx^{k}}
\end{array}\right].\label{eq:MatrixOfDerivatives}
\end{equation}


\subsection{Taylor Series}

Taylor series provide a way in which to approximate infinitely differentiable
functions around a point. The Taylor series of an infinitely differentiable
function $f(x)$ of a single variable around a point $x_{0}$ is

\begin{align}
f\left(x\right) & =f\left(x_{0}\right)+\sum_{i\,\geqslant\,1}\left.\frac{1}{i!}\left(x\,-\,x_{0}\right)^{i}\frac{d^{i}f\left(x\right)}{dx^{i}}\right|_{x\,=\,x_{0}}\nonumber \\
 & =f\left(x_{0}\right)+\sum_{i\,=\,1}^{k}\left.\frac{1}{i!}\left(x\,-\,x_{0}\right)^{i}\frac{d^{i}f\left(x\right)}{dx^{i}}\right|_{x\,=\,x_{0}}+r_{k\,+\,1}(x),\label{eq:UnivariateTaylorSeries}
\end{align}

where $r_{k+1}(x)$ is a remainder term. 

Often, we assume that the remainder term can be ignored, focus on
just the first three terms of \eqref{eq:UnivariateTaylorSeries},
and write

\begin{equation}
f\left(x\right)\approx f\left(x_{0}\right)+\left(x\,-\,x_{0}\right)\left.\frac{df\left(x\right)}{dx}\right|_{x\,=\,x_{0}}+\frac{1}{2}\left(x\,-\,x_{0}\right)^{2}\left.\frac{d^{2}f\left(x\right)}{dx^{2}}\right|_{x\,=\,x_{0}}.\label{eq:TwoTermUnivariateTaylorSeries}
\end{equation}

While a multivariate extension of \eqref{eq:UnivariateTaylorSeries}
exists, it is a three term expansion, analogous to \eqref{eq:TwoTermUnivariateTaylorSeries},
that finds the most use in portfolio construction and risk management:

\begin{equation}
f\left(\mathrm{\mathbf{x}}\right)\approx f\left(\mathrm{\mathbf{x}}_{0}\right)+\left(\mathrm{\mathbf{x}}\,-\,\mathrm{\mathbf{x}}_{0}\right)\left.\nabla f\left(\mathrm{\mathbf{x}}\right)\right|_{\mathrm{\mathbf{x}}=\mathrm{\mathbf{x}}_{0}}+\frac{1}{2}\left(\mathrm{\mathbf{x}}\,-\,\mathrm{\mathbf{x}}_{0}\right)^{\prime}\left.\mathrm{\mathbf{H}}\left(f\left(\mathrm{\mathbf{x}}\right)\right)\right|_{x=x_{0}}\left(\mathrm{\mathbf{x}}\,-\,\mathrm{\mathbf{x}}_{0}\right).\label{eq:TwoTermMultivariateTaylorSeries}
\end{equation}

It is important to keep in mind that a Taylor series expansion of
$f\left(\mathrm{\mathbf{x}}\right)$ is a \emph{local} expansion,
and is typically valid only in a small region centered at $\mathrm{\mathbf{x}}_{0}$.
As $\mathrm{\mathbf{x}}_{0}$ changes, so does its Taylor series,
as well as its gradient and its Hessian. 

\subsection{Convex and Concave Functions\label{subsec:ConvexFunctions} }

A function $f(x)$ of single variable is \emph{concave} if for any
$\alpha\in[0,1]$, and for all $x,y$, 

\begin{equation}
f((1\,-\,\alpha)\,x\,+\,\alpha\,y)\geqslant(1\,-\,\alpha)\,f(x)\,+\,\alpha\,f(y).\label{eq:ConcaveFunction}
\end{equation}
. 

The function $f(x)$ is \emph{strictly concave} if the strict inequality
$f((1\,-\,\alpha)\,x\,+\,\alpha\,y)>(1\,-\,\alpha)\,f(x)\,+\,\alpha\,f(y)$
holds for any $\alpha\in(0,1)$ and $x\ne y$.

A differentiable function $f(x)$ is (strictly) concave on an interval
if and only if its derivative $f^{\prime}(x)$ is (strictly) monotonically
decreasing on that interval. If a function $f$ has a second derivative
$f^{\prime\prime}(x)$ on an interval, then $f$ is (strictly) concave
if and only if $f^{\prime\prime}(x)$ is (negative) non--positive
on that interval.

If $f$ is concave and differentiable, then it is bounded above by
its first--order Taylor approximation:

\begin{equation}
f(x)\leqslant f(x_{0})\,+\,f^{\prime}(x_{0})\,\cdot\,(x\,-\,x_{0})\ \forall x,\,x_{0},\label{eq:ConcaveFunctionUpperBound}
\end{equation}

and has a global \emph{maximum} that is achieved at one of the two
boundaries of the interval or at a point $x_{0}$ for which $f^{\prime}(x_{0})=0.$ 

A function $f(x)$ of single variable is \emph{convex} if $-f(x)$
is concave, i.e. for any $\alpha\in[0,1]$, and for all $x,y$,  $f((1\,-\,\alpha)\,x\,+\,\alpha\,y)\leqslant(1\,-\,\alpha)\,f(x)\,+\,\alpha\,f(y)$.
The function $f(x)$ is \emph{strictly convex} if $f((1\,-\,\alpha)\,x\,+\,\alpha\,y)<(1\,-\,\alpha)\,f(x)\,+\,\alpha\,f(y)\ \forall\alpha\in(0,1)$
and $x\ne y$.

A differentiable function $f(x)$ is (strictly) convex on an interval
if and only if its derivative $f^{\prime}(x)$ is (strictly) monotonically
increasing on that interval. If a function $f$ has a second derivative
$f^{\prime\prime}(x)$ on an interval then $f$ is (strictly) convex
if and only if $f^{\prime\prime}(x)$ is (positive) non--negative
on that interval.

If $f$ is convex and differentiable, it is bounded below by its first-order
Taylor approximation:

\begin{equation}
f(x)\geqslant f(x_{0})\,+\,f^{\prime}(x_{0})\cdot(x\,-\,x_{0})\ \forall\,x,\,x_{0},\label{eq:ConvexFunctionLowerBound}
\end{equation}

and has a global \emph{minimum} that is achieved at one of the two
boundaries of the interval or at a point $x_{0}$ for which $f^{\prime}(x_{0})=0.$ 

\emph{Jensen's inequality}, which sees extensive use in continuous
time finance, concerns the expected value of convex and concave functions
of a random variable, states that for any random variable $Z$, and
for any convex function $f$,

\begin{equation}
E\left[f\left(Z\right)\right]\geqslant f\left(E\left[Z\right]\right).\label{eq:JensensInequalityConvexFunction}
\end{equation}

If $f$ is concave, the direction of the inequality is reversed. 

Jensen's inequality is easily seen to be true by setting $x_{0}=E\left[Z\right]$
in \eqref{eq:ConcaveFunctionUpperBound} and \eqref{eq:ConvexFunctionLowerBound},
and then taking the expectation of both sides.

\subsection{Multivariate Convex and Concave Functions}

To extend the notions of convexity and concavity to functions of many
variables, we need to define a \emph{convex set}. A set $S$ is said
to be convex if for any $\alpha\in[0,1]$, and for all $\mathrm{\mathbf{x}},\mathrm{\mathbf{y}}\in S$,
the line segment joining $\mathrm{\mathbf{x}}$ and $\mathrm{\mathbf{y}}$
lies within, or on the boundary of, $S$, i.e.

\begin{equation}
(1\,-\,\alpha)\,\mathbf{x}\,+\,\alpha\,\mathrm{\mathbf{y}}\in S.\label{Convex set}
\end{equation}

A multivariable function $f\left(\mathrm{\mathbf{x}}\right)$ that
is defined on a convex set $S$ is said to be concave if it satisfies
\eqref{eq:ConcaveFunction} with the scalars $x$ and $y$ being replaced
by vectors $\mathrm{\mathbf{x}}$ and $\mathrm{\mathbf{y}}$. Replacing
scalars with vectors allows us to similarly define strictly concave,
convex and strictly convex functions by analogy to their definitions
in Section \eqref{subsec:ConvexFunctions}. 

A twice--differentiable function $f\left(\mathrm{\mathbf{x}}\right)$
is convex (strictly convex) if its Hessian, $\mathrm{\mathbf{H}}\left(f\left(\mathrm{\mathbf{x}}\right)\right)$,
is positive semi--definite (positive definite), and concave (strictly
concave) if its Hessian is negative semi--definite (negative definite).
Intuitively, a strictly convex function must lie above, and a strictly
concave function must lie below, a hyperplane that is tangent to its
surface at any point $\mathrm{\mathbf{x}}_{0}$. Equation \eqref{eq:TwoTermMultivariateTaylorSeries}
makes clear that this can only be guaranteed for sufficiently small
$\mathrm{\mathbf{x}}\,-\,\mathrm{\mathbf{x}}_{0}$ if the Hessian
is positive definite (for convex functions) or negative definite (for
concave functions).

\subsection{Euler's Theorem for Homogeneous Functions}

A function $f\left(\mathrm{\mathbf{x}}\right)$ is said to be homogeneous
of degree $k>0$ if $f\left(c\mathrm{\mathbf{x}}\right)=c^{k}f\left(\mathrm{\mathbf{x}}\right)\ \forall c>0$. 

Euler’s theorem for homogeneous functions states that $f\left(\mathrm{\mathbf{x}}\right)$
is homogeneous of degree $k$ if and only if 

\begin{equation}
f\left(\mathrm{\mathbf{x}}\right)=\frac{1}{k}\sum_{i\,=\,1}^{n}x_{i}\,\frac{\partial f\left(\mathrm{\mathbf{x}}\right)}{\partial x_{i}}.\label{eq:EulersTheorem1}
\end{equation}

By way of example, $f\left(\mathrm{\mathbf{x}}\right)=x_{1}^{2}\,+\,x_{2}^{2}$
is homogenous of degree 2 because

\begin{align}
f\left(c\,\mathrm{\mathbf{x}}\right) & =c^{2}\left(x_{1}^{2}\,+\,x_{2}^{2}\right)\nonumber \\
 & =c^{2}f\left(\mathrm{\mathbf{x}}\right)\label{eq:EulersTheorem2}
\end{align}

and it is easily seen that

\begin{align}
f\left(\mathrm{\mathbf{x}}\right) & =\frac{1}{2}\left(x_{1}\cdotp2x_{1}\,+\,x_{2}\cdotp2x_{2}\right)\nonumber \\
 & =\frac{1}{2}\left(x_{1}\frac{\partial f\left(\mathrm{\mathbf{x}}\right)}{\partial x_{1}}\,+\,x_{2}\frac{\partial f\left(\mathrm{\mathbf{x}}\right)}{\partial x_{2}}\right).\label{eq:EulersTheorem3}
\end{align}


\section{Optimization}

\subsection{The General Problem of Optimization}

A very general statement of an optimization problem is

\begin{equation}
\begin{aligned}\mathrm{maximize}\ \  & f\left(\mathrm{\mathbf{x}}\right)\\
\mathrm{subject\ to}\\
g\left(\mathrm{\mathbf{x}}\right) & \geqslant0\\
h\left(\mathrm{\mathbf{x}}\right) & =\mathrm{\mathbf{b}}\\
\mathrm{\mathbf{x}} & \in S.
\end{aligned}
\label{eq:GeneralOptimizationProblem}
\end{equation}
This problem is virtually limitless in both its scope and its difficulty:
$f\left(\mathrm{\mathbf{x}}\right),\ g\left(\mathrm{\mathbf{x}}\right)$
and $h\left(\mathrm{\mathbf{x}}\right)$ may be both nonlinear and
discontinuous in their arguments, and $S$ may be bounded, non--convex,
finite or countable (e.g. some or all the components of $\mathrm{\mathbf{x}}$
may be restricted to be integers); $f\left(\mathrm{\mathbf{x}}\right)$
may have multiple local maxima that are not easily distinguished from
its global maximum, and its evaluation may be numerically unstable
and computationally demanding. 

We therefore focus here only on the optimization of sufficiently smooth
(i.e. twice differentiable) concave (or convex) functions defined
on $\mathbb{R}^{n}$ with linear equality constraints. Fortunately,
this covers many of the problems that actually arise in portfolio
management.

\subsection{Optimizing Twice--Differentiable Convex Functions}

Twice--differentiable convex (concave) functions with positive (negative)
definite Hessians play a central role in portfolio construction as
they are easily optimized when $S$, the convex set on which they
are defined, is unbounded. Any extremum of such a function must necessarily
be a minimum or a maximum, depending on whether the Hessian is positive
or negative definite. The sole condition for an extremum is
\begin{equation}
\nabla f\left(\mathrm{\mathbf{x}}_{opt}\right)=\mathbf{0},\label{eq:GradientFxEq0}
\end{equation}

where $\nabla f\left(\mathrm{\mathbf{x}}\right)$ is defined in \eqref{eq:GradientFx}.
If $\mathrm{\mathbf{H}}\left(f\left(\mathrm{\mathbf{x}}\right)\right)$
is positive (negative) definite, $f$ is minimized (maximized) at
$\mathrm{\mathbf{x}}_{opt}$. Intuitively, this follows from \eqref{eq:TwoTermMultivariateTaylorSeries}:
as the gradient is 0 at $\mathrm{\mathbf{x}}_{opt}$, the linear term
disappears, and only the quadratic term survives, so that the function
must increase (decrease) as one moves away from $\mathrm{\mathbf{x}}_{opt}$
in any direction according to whether $\mathrm{\mathbf{H}}\left(f\left(\mathrm{\mathbf{x}}\right)\right)$
is positive (negative) definite.

\subsection{The Method of Lagrange Multipliers}

We next extend our technique for optimizing twice--differentiable
convex (concave) functions with positive (negative) definite Hessians
to permit the inclusion of a single linear equality constraint of
the form

\begin{equation}
\mathrm{\mathbf{a^{\prime}\,x}}=b\mathrm{,}\label{eq:EqualityConstraints}
\end{equation}

where $\mathrm{\mathbf{a}}$ is a vector with constant real entries.

The method of \emph{Lagrange multipliers,} which maps a constrained
optimization to a related unconstrained optimization, is a powerful
technique for solving problems such as this. It starts by constructing
a new function, the \emph{Lagrangian}, that combines the function
and the linear constraint:

\begin{equation}
\mathcal{L}\left(\mathrm{\mathbf{x}},\,\theta\right)=f\left(\mathrm{\mathbf{x}}\right)\,-\,\theta\left(\mathrm{\mathbf{a^{\prime}\,x}}\,-\,b\right).\label{eq:Lagrangian}
\end{equation}

The parameter $\theta$ is known as a Lagrange multiplier, and at
a stationary point of $\mathcal{L}$, we must have 

\begin{equation}
\nabla\mathcal{L}\left(\mathrm{\mathbf{x}},\,\theta\right)=\mathbf{0},\label{eq:GradientLagrangianEq0}
\end{equation}

where

\begin{equation}
\nabla\mathcal{L}\left(\mathrm{\mathbf{x}},\,\theta\right)=\left[\begin{array}{c}
\frac{\partial\mathcal{L}\left(\mathrm{\mathbf{x}},\,\theta\right)}{\partial x_{1}}\\
\vdots\\
\frac{\partial\mathcal{L}\left(\mathrm{\mathbf{x}},\,\theta\right)}{\partial x_{n}}\\
\frac{\partial\mathcal{L}\left(\mathrm{\mathbf{x}},\,\theta\right)}{\partial\theta}
\end{array}\right].\label{eq:GradientLagrangianDefinition}
\end{equation}

The equality constraint \eqref{eq:EqualityConstraints} must be satisfied
at any stationary point of $\mathcal{L}\left(\mathrm{\mathbf{x}},\,\theta\right)$,
as, if it is not, we can increase the value of $\mathcal{L}\left(\mathrm{\mathbf{x}},\,\theta\right)$
by choosing a different value of $\theta$, contradicting our assumption
that it is a stationary point. 

Observe that for any fixed value of $\theta$, the Hessian of $\mathcal{L}\left(\mathrm{\mathbf{x}},\,\theta\right)$
is identical to that of $f\left(\mathrm{\mathbf{x}}\right)$, as differentiating
$\mathcal{L}\left(\mathrm{\mathbf{x}},\,\theta\right)$ twice with
respect to $x_{i}$ and $x_{j}$ for $\theta$ fixed eliminates all
the terms contributed by the linear constraint. In particular, if
the Hessian of $f\left(\mathrm{\mathbf{x}}\right)$ is positive (negative)
definite, so is that of $\mathcal{L}\left(\mathrm{\mathbf{x}},\,\theta\right)$.
It follows that if the Hessian of $f\left(\mathrm{\mathbf{x}}\right)$
is negative definite, $f\left(\mathrm{\mathbf{x}}\right)$ achieves
its constrained maximum at a stationary point of $\mathcal{L}\left(\mathrm{\mathbf{x}},\,\theta\right)$. 

An identical argument allows for the constrained minimization of functions
with positive definite Hessians, and this technique is easily extended
to multiple linear equality constraints -- we need only add one Lagrange
multiplier per constraint, and ensure that the derivative of the Lagrangian
with respect to each Lagrange multiplier is 0.

Once we allow for the inclusion of inequality constraints or non--linear
(typically quadratic) constraints, the optimization problem becomes
much harder, and we refer the reader to specialized textbooks on optimization
such as \citet{LuenbergerYe2008} and \citet{CornuejolsEtAl2018}.
Additionally, \citet{Kwan2007}, \citet{Kwan2018} and \citet{MarkowitzStarerFramGerber2020}
provide useful expositions of optimization models that are of particular
relevance to portfolio management, along with Python and R code, as
well as Excel spreadsheets, to assist in their implementation.

\end{appendices}

\pagebreak{}

\begin{flushleft}
\printbibliography
\par\end{flushleft}

\newpage
\end{document}
